\documentclass[a4paper,11pt]{article}
\usepackage[utf8]{inputenc}
\usepackage[T1]{fontenc}
\usepackage{lmodern}
\usepackage[french]{babel}
\usepackage{amsmath,amssymb}
\usepackage{stmaryrd} % Indispensable pour \llbracket et \rrbracket
\usepackage{xcolor}
\usepackage{tcolorbox}
\usepackage{tikz}
\usetikzlibrary{arrows.meta,positioning}
\tcbuselibrary{skins,breakable}
\usepackage{titlesec}
\usepackage{fancyhdr}
\usepackage[a4paper,margin=2.3cm]{geometry}
\usepackage{microtype}
\usepackage{mathrsfs}

%==================================================================
% COULEURS
%==================================================================
\definecolor{brun}{RGB}{140,90,15}
\definecolor{bordeaux}{RGB}{150,30,45}
\definecolor{bleu}{RGB}{40,100,170}
\definecolor{bleuclair}{RGB}{228,239,250}
\definecolor{vert}{RGB}{30,120,80}
\definecolor{orange}{RGB}{205,100,20}
\definecolor{violet}{RGB}{110,60,150}

%==================================================================
% RACCOURCIS MATHÉMATIQUES
%==================================================================
\newcommand{\K}{\mathbb{K}}
\newcommand{\R}{\mathbb{R}}
\newcommand{\C}{\mathbb{C}}
\newcommand{\N}{\mathbb{N}}
\newcommand{\Z}{\mathbb{Z}}
\newcommand{\Q}{\mathbb{Q}}
\newcommand{\U}{\mathbb{U}}
\newcommand{\D}{\mathcal{D}}
\newcommand{\dd}{\,\mathrm{d}}
\newcommand{\intEnt}[2]{\llbracket #1, #2 \rrbracket}
\DeclareMathOperator{\Argch}{Argch}
\DeclareMathOperator{\sgn}{signe}

% Formule mise en valeur
\newcommand{\cadre}[1]{%
  \tcbox[on line,enhanced,colback=white,colframe=brun,arc=3pt,boxrule=0.7pt,
         boxsep=3pt,left=2pt,right=2pt,top=1pt,bottom=1pt]{$\displaystyle #1$}}

% Mot ou expression mis en valeur dans le texte
\newcommand{\mot}[1]{\textcolor{bordeaux}{\textbf{#1}}}
\newcommand{\cas}[1]{\textbf{\textcolor{orange}{#1}}}

%==================================================================
% PAGE, EN-TÊTE, TITRES (Aéré et espacé)
%==================================================================
\pagestyle{fancy}
\fancyhf{}
\fancyhead[L]{\small\textcolor{brun}{Mathématiques}}
\fancyhead[R]{\small\textcolor{brun}{Chapitre 12 -- Lois de composition internes / Groupes}}
\fancyfoot[C]{\small\thepage}
\renewcommand{\headrule}{\hbox to\headwidth{\color{brun}\leaders\hrule height \headrulewidth\hfill}}
\setlength{\parindent}{0pt}
\setlength{\parskip}{8pt}
\renewcommand{\labelitemi}{\textcolor{bleu}{\textbullet}}

\titleformat{\section}{\Large\scshape\color{bordeaux}}{}{0pt}{}[{\color{bordeaux}\titlerule[0.8pt]}]
\titleformat{\subsection}{\large\scshape\color{brun}}{}{0pt}{}
\titlespacing*{\section}{0pt}{20pt}{10pt}
\titlespacing*{\subsection}{0pt}{16pt}{8pt}

%==================================================================
% BLOCS COLORÉS
%==================================================================
\newcommand{\titrebloc}[2]{\textsc{#1}\ifx\relax#2\relax\else\ \textsc{(#2)}\fi}
\newcommand{\titreblocnum}[2]{\textsc{#1}\ifx\relax#2\relax\else~#2\fi}

\tcbset{
  bloc/.style={enhanced,breakable,arc=3pt,boxrule=0.8pt,
    left=8pt,right=8pt,top=7pt,bottom=7pt,
    fonttitle=\small,attach title to upper={\ \ },
    before skip=14pt,after skip=14pt},
  marge/.style={enhanced,breakable,frame hidden,arc=2pt,
    left=10pt,right=6pt,top=6pt,bottom=6pt,
    fonttitle=\small,attach title to upper={\ \ },
    before skip=12pt,after skip=12pt}
}

% Environnements de base
\newtcolorbox{definition}[1][]{bloc, colback=brun!7,colframe=brun,coltitle=brun,colbacktitle=brun!7, title={\titrebloc{Définition}{#1}}}
\newtcolorbox{proposition}[1][]{bloc, colback=bordeaux!5,colframe=bordeaux,coltitle=bordeaux,colbacktitle=bordeaux!5, title={\titrebloc{Proposition}{#1}}}
\newtcolorbox{theoreme}[1][]{bloc, colback=bordeaux!8,colframe=bordeaux,boxrule=1.4pt, coltitle=bordeaux,colbacktitle=bordeaux!8, title={\titrebloc{Théorème}{#1}}}
\newtcolorbox{demonstration}{marge, colback=bleuclair,borderline west={2.5pt}{0pt}{bleu}, before upper={\textcolor{bleu}{\small\scshape Démonstration.}\ }, after upper={\par\nopagebreak\hfill\textcolor{bleu}{$\blacksquare$}}}
\newtcolorbox{remarque}[1][]{marge, colback=vert!6,borderline west={2.5pt}{0pt}{vert}, coltitle=vert!70!black,colbacktitle=vert!6, title={\titreblocnum{Remarque}{#1}}}
\newtcolorbox{exemples}[1][Exemples]{marge, colback=orange!6,borderline west={2.5pt}{0pt}{orange}, coltitle=orange!85!black,colbacktitle=orange!6, title={\textsc{#1}}}
\newtcolorbox{methode}[1][Méthode]{marge, colback=violet!6,borderline west={2.5pt}{0pt}{violet}, coltitle=violet,colbacktitle=violet!6, title={\textsc{#1}}}

% Nouveaux environnements
\newtcolorbox{lemme}[1][]{bloc, colback=bleu!5,colframe=bleu,coltitle=bleu,colbacktitle=bleu!5, title={\titrebloc{Lemme}{#1}}}
\newtcolorbox{corollaire}[1][]{bloc, colback=orange!5,colframe=orange,coltitle=orange,colbacktitle=orange!5, title={\titrebloc{Corollaire}{#1}}}
\newtcolorbox{synthese}[1][]{bloc, colback=violet!5,colframe=violet,coltitle=violet,colbacktitle=violet!5, title={\titrebloc{Synthèse}{#1}}}

%==================================================================
% DEBUT DU DOCUMENT
%==================================================================
\begin{document}

\begin{center}
  {\Huge \scshape \color{bordeaux} Chapitre 12 -- Lois de composition internes / Groupes}
\end{center}
\vspace{10pt}

\section{I -- Loi de composition interne}

\subsection{1.1. Définitions}

\textbf{a)}
\begin{definition}
On appelle \mot{loi de composition interne} sur un ensemble $E$ toute application
\[
T : E \times E \longrightarrow E
\]
de $E \times E$ dans $E$. ($T$ se lit \textit{« truc »}).

Le couple $(E, T)$ s'appelle un \mot{magma}.
\end{definition}

\begin{remarque}
Soit $T : E \times E \to E$ une loi de composition interne sur $E$. Il est d'usage de noter l'image d'un couple $(x, y) \in E^2$ par $x \, T \, y$ de préférence à $T(x, y)$.

Souvent, la loi de composition $T$ est notée par un point $x \cdot y$, voire sans symbole $xy$ (\mot{notation multiplicative}).

Dans ce cas, si $(x, y, z) \in E^3$, $x(yz)$ désigne $x \, T \, (y \, T \, z)$.
\end{remarque}

\textbf{b)} Soient $x_1, \dots, x_n$ $n$ éléments d'un magma $(E, T)$. Pour donner un sens à $x_1 \, T \, x_2 \, T \dots T \, x_n$, il faut préciser le \mot{parenthésage}.

\begin{exemples}[Exemples]
\begin{itemize}
    \item Pour $n = 3$ :
    \[
    x_1 \, T \, (x_2 \, T \, x_3), \qquad (x_1 \, T \, x_2) \, T \, x_3
    \]
    \item Pour $n = 4$ :
    \[
    (x_1 \, T \, x_2) \, T \, (x_3 \, T \, x_4), \qquad x_1 \, T \, (x_2 \, T \, (x_3 \, T \, x_4)), \qquad ((x_1 \, T \, x_2) \, T \, x_3) \, T \, x_4,
    \]
    \[
    (x_1 \, T \, (x_2 \, T \, x_3)) \, T \, x_4, \qquad x_1 \, T \, ((x_2 \, T \, x_3) \, T \, x_4)
    \]
\end{itemize}
\end{exemples}

Il est commode d'encoder un parenthésage sur $n$ symboles par des \mot{arbres binaires planaires enracinés à $n$ feuilles} (chaque nœud interne possède exactement $2$ fils, avec un seul sommet sans parent dit la \mot{racine}, et les nœuds sans fils appelés \mot{feuilles} numérotés de gauche à droite).

\begin{exemples}[Représentation par arbres binaires planaires]
\textbf{Pour $n = 3$ :}
\begin{center}
\begin{tikzpicture}[
  level 1/.style={sibling distance=16mm, level distance=10mm},
  level 2/.style={sibling distance=11mm, level distance=10mm},
  noeud/.style={circle, fill=bleu, inner sep=1.6pt},
  feuille/.style={font=\small}
]
  % Arbre 1 : (x1 T x2) T x3
  \begin{scope}[xshift=-3.2cm]
    \node[noeud] (r1) {}
      child { node[noeud] {}
        child { node[feuille] {$x_1$} }
        child { node[feuille] {$x_2$} }
      }
      child { node[feuille] {$x_3$} };
    \node[above=2pt of r1, font=\scriptsize\color{brun}] {racine};
    \node at (0, -2.6) {$(x_1 \, T \, x_2) \, T \, x_3$};
  \end{scope}

  \node at (0, -1.0) {\Large $\neq$};

  % Arbre 2 : x1 T (x2 T x3)
  \begin{scope}[xshift=3.2cm]
    \node[noeud] (r2) {}
      child { node[feuille] {$x_1$} }
      child { node[noeud] {}
        child { node[feuille] {$x_2$} }
        child { node[feuille] {$x_3$} }
      };
    \node at (0, -2.6) {$x_1 \, T \, (x_2 \, T \, x_3)$};
  \end{scope}
\end{tikzpicture}
\end{center}
\textit{(Les deux arbres sont distincts en tant qu'arbres planaires, alors que dans l'espace ils seraient isomorphes).}

\vspace{12pt}
\textbf{Pour $n = 4$ (Graphe des transformations élémentaires) :}
\begin{center}
\begin{tikzpicture}[
  level 1/.style={sibling distance=16mm, level distance=9mm},
  level 2/.style={sibling distance=10mm, level distance=9mm},
  level 3/.style={sibling distance=9mm, level distance=9mm},
  noeud/.style={circle, fill=bleu, inner sep=1.4pt},
  feuille/.style={font=\scriptsize}
]
  % Sommet haut : (x1 T x2) T (x3 T x4)
  \begin{scope}[yshift=3.8cm]
    \node[noeud] {}
      child { node[noeud] {}
        child { node[feuille] {$x_1$} }
        child { node[feuille] {$x_2$} }
      }
      child { node[noeud] {}
        child { node[feuille] {$x_3$} }
        child { node[feuille] {$x_4$} }
      };
    \node at (0, -2.3) {$(x_1 \, T \, x_2) \, T \, (x_3 \, T \, x_4)$};
  \end{scope}

  % Sommet gauche : x1 T (x2 T (x3 T x4))
  \begin{scope}[xshift=-4.8cm, yshift=0.5cm]
    \node[noeud] {}
      child { node[feuille] {$x_1$} }
      child { node[noeud] {}
        child { node[feuille] {$x_2$} }
        child { node[noeud] {}
          child { node[feuille] {$x_3$} }
          child { node[feuille] {$x_4$} }
        }
      };
    \node at (0, -3.1) {$x_1 \, T \, (x_2 \, T \, (x_3 \, T \, x_4))$};
  \end{scope}

  % Sommet droite : ((x1 T x2) T x3) T x4
  \begin{scope}[xshift=4.8cm, yshift=0.5cm]
    \node[noeud] {}
      child { node[noeud] {}
        child { node[noeud] {}
          child { node[feuille] {$x_1$} }
          child { node[feuille] {$x_2$} }
        }
        child { node[feuille] {$x_3$} }
      }
      child { node[feuille] {$x_4$} };
    \node at (0, -3.1) {$((x_1 \, T \, x_2) \, T \, x_3) \, T \, x_4$};
  \end{scope}

  % Sommet bas : (x1 T (x2 T x3)) T x4
  \begin{scope}[yshift=-2.8cm]
    \node[noeud] {}
      child { node[noeud] {}
        child { node[feuille] {$x_1$} }
        child { node[noeud] {}
          child { node[feuille] {$x_2$} }
          child { node[feuille] {$x_3$} }
        }
      }
      child { node[feuille] {$x_4$} };
    \node at (0, -3.1) {$(x_1 \, T \, (x_2 \, T \, x_3)) \, T \, x_4$};
  \end{scope}

  % Arêtes des transformations élémentaires
  \draw[<->, >=stealth, thick, bordeaux] (-1.5, 2.7) -- (-3.6, 0.9);
  \draw[<->, >=stealth, thick, bordeaux] (1.5, 2.7) -- (3.6, 0.9);
  \draw[<->, >=stealth, thick, bordeaux] (-3.2, -2.9) -- (-1.3, -4.1);
  \draw[<->, >=stealth, thick, bordeaux] (3.2, -2.9) -- (1.3, -4.1);
\end{tikzpicture}
\end{center}
\end{exemples}

On peut montrer que le nombre de parenthésages différents sur $n+1$ éléments $x_1, \dots, x_{n+1}$ de $E$ est :
\[
\cadre{C_n = \frac{1}{n+1} \binom{2n}{n}}
\]
appelé \mot{nombre de Catalan}.

On peut démontrer que l'on peut passer d'un parenthésage à un autre sur $x_1, \dots, x_n$ (dans cet ordre) en appliquant une succession finie de \mot{transformations élémentaires} :
\[
(xy)z \longrightarrow x(yz) \quad \text{ou} \quad x(yz) \longrightarrow (xy)z
\]
comme on peut le vérifier pour $n = 4$, où les transformations élémentaires sont représentées par les arêtes du graphe ci-dessus.

\vspace{8pt}
\textbf{c)} À tout magma $(E, T)$, on peut associer le magma $(\mathcal{P}(E), *)$ en définissant $*$ par :
\[
\forall (A, B) \in \mathcal{P}(E)^2, \quad A * B = \{ z \in E \mid \exists (x, y) \in A \times B, \ z = x \, T \, y \}
\]

Par abus de notation, on note $A \, T \, B$ pour $A * B$, et même $x \, T \, B$ pour $\{x\} * B$ et $A \, T \, y$ pour $A * \{y\}$.

\begin{exemples}
Dans $\R$, somme de deux parties $A$ et $B$ :
\[
A + B = \{ z \in \R \mid \exists (x, y) \in A \times B, \ z = x + y \}
\]
En particulier, pour $x \in \R$ :
\[
x + 2\pi\Z = \{ z \in \R \mid \exists k \in \Z, \ z = x + 2k\pi \}
\]
\end{exemples}

\textbf{d)}
\begin{definition}
Soit $(E, T)$ un magma et $a \in E$. On appelle \mot{translation à gauche} (resp. \mot{à droite}) par $a$ l'application :
\[
\gamma_a : E \longrightarrow E, \quad x \longmapsto a \, T \, x \qquad \text{(resp. } \delta_a : E \longrightarrow E, \quad x \longmapsto x \, T \, a \text{)}
\]
\end{definition}

\begin{definition}[Partie stable -- Loi induite]
Soit $(E, T)$ un magma et $A$ une partie de $E$. On dit que $A$ est \mot{stable} pour $T$ si :
\[
\forall (x, y) \in A^2, \quad x \, T \, y \in A
\]
L'application
\[
T_A : A \times A \longrightarrow A, \quad (x, y) \longmapsto x \, T \, y
\]
est alors une loi de composition interne sur $A$, appelée \mot{loi induite} par $T$ sur $A$.

S'il n'y a pas d'ambiguïté, on note encore $T$ pour $T_A$.
\end{definition}

\newpage
\subsection{1.2. Lois associatives}

\begin{definition}
Soit $(E, T)$ un magma. La loi $T$ est dite \mot{associative} si :
\[
\forall (x, y, z) \in E^3, \quad x \, T \, (y \, T \, z) = (x \, T \, y) \, T \, z
\]
On dit alors que $(E, T)$ est un \mot{magma associatif}.
\end{definition}

Si $P$ désigne un parenthésage sur $n$ symboles, et $x_1, \dots, x_n$ sont $n$ éléments de $E$, on note $P(x_1, \dots, x_n)$ l'élément de $E$ obtenu en composant $x_1, \dots, x_n$ par la loi $T$ en suivant le parenthésage $P$.

\begin{exemples}
\begin{itemize}
    \item Pour $n = 3$ :
    \begin{align*}
        P = ((\bullet, \bullet), \bullet) &\implies P(x, y, z) = (x \, T \, y) \, T \, z \\
        P = (\bullet, (\bullet, \bullet)) &\implies P(x, y, z) = x \, T \, (y \, T \, z)
    \end{align*}
    \item Pour $n = 4$ :
    \begin{align*}
        P = ((\bullet, \bullet), (\bullet, \bullet)) &\implies P(x, y, z, t) = (x \, T \, y) \, T \, (z \, T \, t) \\
        P = (\bullet, (\bullet, (\bullet, \bullet))) &\implies P(x, y, z, t) = x \, T \, (y \, T \, (z \, T \, t))
    \end{align*}
\end{itemize}
\end{exemples}

\begin{proposition}
Soit $(E, T)$ un magma associatif et $x_1, \dots, x_n$ $n$ éléments de $E$. Si $P$ et $P'$ sont deux parenthésages sur $n$ symboles, on a :
\[
P(x_1, \dots, x_n) = P'(x_1, \dots, x_n)
\]
En d'autres termes, on peut noter $x_1 \, T \, x_2 \, T \dots T \, x_n$ sans préciser le parenthésage.
\end{proposition}

\begin{demonstration}
\textit{(Admise)} Résulte de la remarque en fin de \textbf{b)} sur les transformations élémentaires : dans un magma associatif, chaque transformation élémentaire conserve la valeur du produit puisque $(xy)z = x(yz)$.
\end{demonstration}

\begin{remarque}[Notation]
Lorsque $x_1 = \dots = x_n = x$ ($n \in \N^*$) dans un magma associatif, on note :
\[
\overset{n}{\underset{}{\top}} x = \underbrace{x \, T \, x \, T \dots T \, x}_{n \text{ fois}}
\]
En notation multiplicative, on note cela $x^n$.
\end{remarque}

\newpage
\subsection{1.3. Lois commutatives}

\begin{definition}
Soit $(E, T)$ un magma.
\begin{enumerate}
    \item On dit que deux éléments $a$ et $b$ de $E$ sont \mot{permutables}, ou \mot{commutent}, si :
    \[
    a \, T \, b = b \, T \, a
    \]
    \item La loi $T$ est dite \mot{commutative} si deux éléments quelconques de $E$ sont permutables, c'est-à-dire :
    \[
    \forall (a, b) \in E^2, \quad a \, T \, b = b \, T \, a
    \]
    On dit alors que $(E, T)$ est un \mot{magma commutatif}.
\end{enumerate}
\end{definition}

\begin{proposition}
Soit $(E, T)$ un magma associatif, et $x_1, \dots, x_n$ des éléments de $E$ deux à deux permutables (sans avoir nécessairement la commutativité sur tout $E$). Le composé $x_1 \, T \, x_2 \, T \dots T \, x_n$ de ces éléments ne dépend pas de l'ordre dans lequel on les compose.
\end{proposition}

\begin{demonstration}
Il s'agit de prouver que pour toute permutation $\sigma : \intEnt{1}{n} \to \intEnt{1}{n}$, on a :
\[
x_{\sigma(1)} \, T \, x_{\sigma(2)} \, T \dots T \, x_{\sigma(n)} = x_1 \, T \, x_2 \, T \dots T \, x_n
\]

Cette égalité est vraie si $\sigma$ est une \mot{transposition élémentaire} $\tau_i$ (où $1 \le i < n$), définie par :
\[
\tau_i(i) = i+1, \qquad \tau_i(i+1) = i, \qquad \text{et} \quad \forall j \notin \{i, i+1\}, \quad \tau_i(j) = j
\]
En effet, puisque $x_i$ et $x_{i+1}$ commutent et que $T$ est associative :
\begin{align*}
    x_1 \, T \, x_2 \, T \dots T \, x_{\tau_i(i)} \, T \, x_{\tau_i(i+1)} \, T \dots T \, x_n 
    &= x_1 \, T \, x_2 \, T \dots T \, (x_{i+1} \, T \, x_i) \, T \dots T \, x_n \\
    &= x_1 \, T \, x_2 \, T \dots T \, (x_i \, T \, x_{i+1}) \, T \dots T \, x_n \\
    &= x_1 \, T \, x_2 \, T \dots T \, x_i \, T \, x_{i+1} \, T \dots T \, x_n
\end{align*}

Toute permutation de $\mathfrak{S}_n$ étant composée de transpositions élémentaires, l'égalité est vraie pour tout $\sigma \in \mathfrak{S}_n$.
\end{demonstration}

\begin{remarque}[Conséquences]
\begin{enumerate}
    \item Le composé pourra être noté $\displaystyle\underset{i=1}{\overset{n}{\top}} x_i$ sans que l'ordre dans lequel on compose les $x_i$ soit imposé.
    \item Soient $x, y$ deux éléments permutables pour une loi associative $T$. Alors :
    \begin{itemize}
        \item $\forall n, p \in \N^*$, $\overset{n}{\top} x$ et $\overset{p}{\top} y$ sont permutables (en notation multiplicative, $x^n$ et $y^p$ sont permutables).
        \item En particulier, $\forall n \in \N^*$, $(\overset{n}{\top} x) \, T \, (\overset{n}{\top} y) = \overset{n}{\top} (x \, T \, y)$ (en notation multiplicative, $x^n \cdot y^n = (x \cdot y)^n$).
    \end{itemize}
    \textit{(Preuve : par récurrence sur $n$, exercice à faire).}
\end{enumerate}
\end{remarque}

\newpage
\subsection{1.4. Élément neutre}

\begin{definition}
Soit $(E, T)$ un magma. On appelle \mot{élément neutre} pour la loi $T$ tout élément $e$ de $E$ tel que :
\[
\forall x \in E, \quad e \, T \, x = x \, T \, e = x
\]
\end{definition}

\begin{proposition}
S'il existe, un tel élément neutre est unique.
\end{proposition}

\begin{demonstration}
Soient $e$ et $e' \in E$ deux éléments neutres. Alors :
\[
e \, T \, e' = e' \quad (\text{car } e \text{ est neutre})
\]
\[
e \, T \, e' = e \quad (\text{car } e' \text{ est neutre})
\]
Donc $e = e'$.
\end{demonstration}

\begin{definition}[Magma unifère et Monoïde]
\begin{itemize}
    \item Un magma $(E, T)$ qui possède un élément neutre est dit \mot{unifère}.
    \item On appelle \mot{monoïde} tout magma associatif et unifère $(E, T)$. Lorsque de plus $T$ est commutative, le monoïde est dit \mot{commutatif} ou \mot{abélien}.
\end{itemize}
\end{definition}

\begin{remarque}[Convention]
Si $(E, T)$ est un magma associatif et unifère de neutre $e$, on étend la notation $\overset{n}{\top} x$ à $n = 0$ en posant :
\[
\overset{0}{\top} x = e \qquad \text{(en notation multiplicative : } x^0 = e\text{)}
\]
\end{remarque}

\begin{exemples}
\begin{enumerate}
    \item $(\N, +)$ est un monoïde commutatif dont l'élément neutre est $0$.
    \item $(\N, \times)$ est un monoïde commutatif dont l'élément neutre est $1$.
    \item Soit $X$ un ensemble. Le magma $(X^X, \circ)$ est un monoïde dont l'élément neutre est $\mathrm{id}_X$, muni de la loi :
    \[
    X^X \times X^X \longrightarrow X^X, \quad (f, g) \longmapsto f \circ g
    \]
    En effet, la composition des applications est associative :
    \[
    \forall f, g, h \in X^X, \quad (f \circ g) \circ h = f \circ (g \circ h)
    \]
    et possède un neutre :
    \[
    \forall f \in X^X, \quad f \circ \mathrm{id}_X = \mathrm{id}_X \circ f = f
    \]
    \item $(\mathcal{P}(X), \cup)$ et $(\mathcal{P}(X), \cap)$ sont des monoïdes commutatifs, de neutres respectifs $\emptyset$ et $X$.
\end{enumerate}
\end{exemples}

\newpage
\subsection{1.5. Éléments réguliers (ou simplifiables) -- Éléments symétrisables (ou inversibles)}

\begin{definition}[Élément régulier]
Soit $(E, T)$ un magma. On appelle \mot{élément régulier} (ou \mot{simplifiable}) \mot{à droite} (resp. \mot{à gauche}) tout élément $a \in E$ tel que :
\[
\forall (x, y) \in E^2, \quad x \, T \, a = y \, T \, a \implies x = y
\]
\[
\text{(resp. } \forall (x, y) \in E^2, \quad a \, T \, x = a \, T \, y \implies x = y \text{)}
\]
c'est-à-dire tel que $\delta_a$ (resp. $\gamma_a$) soit injective.

Un élément est dit \mot{régulier} (ou \mot{simplifiable}) s'il est régulier à droite et à gauche.
\end{definition}

\begin{exemples}
\begin{enumerate}
    \item Tout élément de $(\N, +)$ est simplifiable :
    \[
    \forall m, n, p \in \N, \quad m + p = n + p \implies m = n
    \]
    \item Tout élément de $(\N^*, \times)$ est simplifiable. En revanche, $(\N, \times)$ possède un seul élément non simplifiable : $0$ (par exemple, $0 \times 0 = 0 \times 1$, mais $0 = 1$ est faux).
    \item Dans le monoïde $(X^X, \circ)$ :
    \begin{itemize}
        \item Les éléments $f \in X^X$ réguliers à droite, c'est-à-dire tels que :
        \[
        \forall g, h \in X^X, \quad g \circ f = h \circ f \implies g = h
        \]
        sont exactement les \mot{surjections}.
        \item Les éléments $f \in X^X$ réguliers à gauche, c'est-à-dire tels que :
        \[
        \forall g, h \in X^X, \quad f \circ g = f \circ h \implies g = h
        \]
        sont exactement les \mot{injections}.
    \end{itemize}
    \item En général, les monoïdes $(\mathcal{P}(X), \cup)$ et $(\mathcal{P}(X), \cap)$ ne sont pas réguliers (dès que $X \neq \emptyset$).
\end{enumerate}
\end{exemples}

\begin{definition}[Élément symétrisable]
Soit $(E, T)$ un magma unifère de neutre $e$. On appelle \mot{élément symétrisable} (resp. \mot{symétrisable à droite}, resp. \mot{à gauche}) tout élément $a \in E$ tel qu'il existe au moins un élément $a' \in E$ tel que :
\[
a \, T \, a' = a' \, T \, a = e \qquad \text{(resp. } a \, T \, a' = e, \quad \text{resp. } a' \, T \, a = e \text{)}
\]
On dit alors que $a'$ est un \mot{symétrique} ou \mot{inverse} (resp. \mot{à droite}, resp. \mot{à gauche}) de $a$ pour la loi $T$.
\end{definition}

\begin{remarque}
\textbf{Attention à l'unicité :} dans un magma unifère quelconque, un symétrique à gauche ou à droite n'est pas forcément unique.
\end{remarque}

\begin{theoreme}
Soit $(E, T)$ un monoïde d'élément neutre $e$, et $a \in E$. Alors :
\begin{enumerate}
    \item Si $a$ est symétrisable à droite (resp. à gauche), alors $a$ est régulier à droite (resp. à gauche).
    \item Si $a$ admet un symétrique à droite et un symétrique à gauche, ceux-ci sont égaux.
    \item Si $a$ est symétrisable, il admet un symétrique unique.
    \item Le composé de deux éléments symétrisables est symétrisable.
\end{enumerate}
\end{theoreme}

\begin{demonstration}
\begin{enumerate}
    \item Supposons $a$ symétrisable à droite, donc qu'il existe $a' \in E$ tel que $a \, T \, a' = e$.
    Soient $x, y \in E$ tels que $x \, T \, a = y \, T \, a$. Alors :
    \[
    (x \, T \, a) \, T \, a' = (y \, T \, a) \, T \, a'
    \]
    Donc, par associativité de $T$ :
    \[
    x \, T \, (a \, T \, a') = y \, T \, (a \, T \, a') \implies x \, T \, e = y \, T \, e \implies x = y
    \]
    Donc $a$ est régulier à droite. De même si $a$ est symétrisable à gauche.

    \item Supposons qu'il existe $a'$ et $a'' \in E$ tels que $a \, T \, a' = e$ et $a'' \, T \, a = e$. Par associativité de $T$ :
    \[
    (a'' \, T \, a) \, T \, a' = a'' \, T \, (a \, T \, a') \implies e \, T \, a' = a'' \, T \, e \implies a' = a''
    \]

    \item Si $a'$ et $a''$ sont deux symétriques de $a$, alors en particulier $a'$ est symétrique à droite et $a''$ est symétrique à gauche, d'où $a' = a''$ d'après le point 2).

    \item Si $a$ et $b$ sont symétrisables de symétriques respectifs $a'$ et $b'$, alors :
    \begin{align*}
        (a \, T \, b) \, T \, (b' \, T \, a') &= a \, T \, (b \, T \, b') \, T \, a' \\
                                              &= a \, T \, e \, T \, a' \\
                                              &= a \, T \, a' = e
    \end{align*}
    et de même :
    \begin{align*}
        (b' \, T \, a') \, T \, (a \, T \, b) &= b' \, T \, (a' \, T \, a) \, T \, b \\
                                              &= b' \, T \, e \, T \, b \\
                                              &= b' \, T \, b = e
    \end{align*}
    Donc $a \, T \, b$ est symétrisable de symétrique $b' \, T \, a'$.
\end{enumerate}
\end{demonstration}

\begin{exemples}
\begin{enumerate}
    \item Le neutre $e$ d'un monoïde est toujours symétrisable, égal à son symétrique ($e \, T \, e = e$).
    \item Dans le monoïde $(\mathcal{P}(X), \cup)$ (de neutre $\emptyset$), le seul élément $A \in \mathcal{P}(X)$ symétrisable (c'est-à-dire tel qu'il existe $B \in \mathcal{P}(X)$ vérifiant $A \cup B = \emptyset$) est l'ensemble $\emptyset$.
    
    Dans le monoïde $(\mathcal{P}(X), \cap)$ (de neutre $X$), le seul élément symétrisable est $X$.
    \item Dans $(\N, +)$, de neutre $0$, le seul élément symétrisable est $0$, mais tous les éléments sont réguliers.
    \item Dans le monoïde $(X^X, \circ)$ :
    \begin{itemize}
        \item[\textbf{a)}] Les éléments $f \in X^X$ symétrisables à droite, c'est-à-dire tels qu'il existe $g \in X^X$ vérifiant $f \circ g = \mathrm{id}_X$, sont exactement les \mot{surjections} (ce qui coïncide ici avec les éléments réguliers à droite).
        \item[\textbf{b)}] Les éléments $f \in X^X$ symétrisables à gauche, c'est-à-dire tels qu'il existe $g \in X^X$ vérifiant $g \circ f = \mathrm{id}_X$, sont exactement les \mot{injections} (ce qui coïncide ici avec les éléments réguliers à gauche).
    \end{itemize}
\end{enumerate}
\end{exemples}

\begin{remarque}[Notation $\overset{n}{\top} x$ étendue à $\Z$]
Soit $(E, T)$ un monoïde de neutre $e$. Si $x \in E$ est symétrisable de symétrique $x'$, on étend la notation $\overset{n}{\top} x$ à $n \in \Z$ en posant :
\[
\forall n \in \Z \setminus \N, \quad \overset{n}{\top} x = \overset{-n}{\top} x'
\]
En notation multiplicative : $x^n = (x')^{-n}$ (on peut donc noter $x^{-1} = x'$).
\end{remarque}

\newpage
\section{II -- Structure de groupe}

\subsection{2.1. Définitions, premières propriétés}

\begin{definition}
On appelle \mot{groupe} tout magma $(G, T)$ vérifiant les propriétés suivantes :
\begin{itemize}
    \item \textbf{(A)} $\forall x, y, z \in G, \quad (x \, T \, y) \, T \, z = x \, T \, (y \, T \, z)$ \quad (associativité)
    \item \textbf{(N)} $\exists e \in G, \ \forall x \in G, \quad x \, T \, e = e \, T \, x = x$ \quad (existence d'un élément neutre)
    \item \textbf{(S)} $\forall x \in G, \ \exists x' \in G, \quad x \, T \, x' = x' \, T \, x = e$ \quad (tout élément est symétrisable)
\end{itemize}
En d'autres termes, un groupe est un \mot{monoïde dont tout élément est symétrisable}.
On dit que $G$ est l'\mot{ensemble sous-jacent} du groupe $(G, T)$.
\end{definition}

Dans la pratique, $G$ désigne aussi bien le groupe $(G, T)$ que l'ensemble sous-jacent.

\begin{remarque}
Un groupe n'est \mot{jamais vide} (il contient au moins l'élément neutre $e$).
\end{remarque}

\begin{remarque}[Notations]
\begin{itemize}
    \item Les lois de groupes sont le plus souvent notées \mot{multiplicativement} $(G, \cdot)$. Le neutre $e$ de $G$ est unique (il peut se noter $1$ ou $1_G$). Si $x \in G$, le symétrique de $x$ (qui est unique) est noté $x'$ ou $x^{-1}$.
    \item Lorsque la loi est commutative (on dit alors que le groupe est \mot{commutatif} ou \mot{abélien}), elle pourra être notée \mot{additivement} $(G, +)$. Le neutre se note alors $0$ (ou $0_G$) et le symétrique de $x$ se note $-x$ (on l'appelle également l'\mot{opposé} de $x$).
\end{itemize}
\end{remarque}

\begin{exemples}
\begin{enumerate}
    \item $(\Z, +)$ est un groupe abélien. En revanche, $(\Z, \times)$ est un monoïde (de neutre $1$), mais seuls $1$ et $-1$ sont symétrisables.
    
    $(\Q, +)$ et $(\Q^*, \times)$ sont deux groupes commutatifs :
    \[
    \forall q \in \Q^*, \ \exists! q' \in \Q^*, \quad qq' = 1
    \]
    mais $0$ n'est pas symétrisable dans le monoïde $(\Q, \times)$ ($\forall q \in \Q, \ q \times 0 = 0 \neq 1$).
    
    De même, $(\R, +)$, $(\R^*, \times)$, $(\C, +)$, $(\C^*, \times)$ sont des groupes commutatifs.
    
    \item $(\N, +)$ n'est pas un groupe, puisque seul $0$ est symétrisable.
    
    \item Si $E$ est un ensemble, $(\mathfrak{S}_E, \circ)$, le monoïde des permutations de $E$, est un groupe, appelé le \mot{groupe des permutations} de $E$ (non abélien dès que $\mathrm{Card}(E) \ge 3$).
    
    En effet, la loi $\mathfrak{S}_E \times \mathfrak{S}_E \to \mathfrak{S}_E, \ (\sigma, \tau) \mapsto \sigma \circ \tau$ vérifie :
    \begin{itemize}
        \item \textbf{(A)} $\forall \sigma_1, \sigma_2, \sigma_3 \in \mathfrak{S}_E, \quad (\sigma_1 \circ \sigma_2) \circ \sigma_3 = \sigma_1 \circ (\sigma_2 \circ \sigma_3)$
        \item \textbf{(N)} $\forall \sigma \in \mathfrak{S}_E, \quad \sigma \circ \mathrm{id}_E = \mathrm{id}_E \circ \sigma = \sigma$
        \item \textbf{(S)} $\forall \sigma \in \mathfrak{S}_E, \quad \sigma \circ \sigma^{-1} = \sigma^{-1} \circ \sigma = \mathrm{id}_E$ (où $\sigma^{-1} \in \mathfrak{S}_E$ est la bijection réciproque).
    \end{itemize}
    En particulier, si $E = \intEnt{1}{n}$ où $n \in \N^*$, $\mathfrak{S}_E$ se note $\mathfrak{S}_n$, appelé \mot{groupe symétrique} d'indice $n$.
\end{enumerate}
\end{exemples}

\begin{proposition}[Groupe des inversibles d'un monoïde]
Plus généralement, si $(M, \cdot)$ est un monoïde, l'ensemble $\mathcal{U}_M$ de ses éléments symétrisables forme un groupe pour la loi induite, appelé \mot{groupe des inversibles} du monoïde $(M, \cdot)$.
\end{proposition}

\begin{demonstration}
\begin{itemize}
    \item \textbf{Stabilité :} Si $x, y \in \mathcal{U}_M$, alors on sait que dans le monoïde $(M, \cdot)$, $x \cdot y$ est symétrisable (de symétrique $y^{-1} x^{-1}$). Donc la loi $\cdot$ de $(M, \cdot)$ induit une loi interne sur $\mathcal{U}_M$.
    \item \textbf{(A)} L'associativité de la loi induite sur $\mathcal{U}_M$ découle de celle de $(M, \cdot)$.
    \item \textbf{(N)} On sait que $e \in \mathcal{U}_M$ (car $e \cdot e = e$) et $e$ est encore neutre pour la loi induite sur $\mathcal{U}_M$.
    \item \textbf{(S)} Soit $x \in \mathcal{U}_M$. Alors il existe $x' \in M$ tel que $x \cdot x' = x' \cdot x = e$. Donc $x'$ est symétrisable de symétrique $x$, d'où $x' \in \mathcal{U}_M$.
\end{itemize}
En conclusion, $(\mathcal{U}_M, \cdot)$ est un groupe.
\end{demonstration}

\begin{exemples}[Suite des exemples de groupes]
Pour $(M, \cdot) = (E^E, \circ)$, les éléments symétrisables $f \in E^E$ de ce monoïde sont les fonctions $f : E \to E$ telles qu'il existe $g : E \to E$ vérifiant $f \circ g = g \circ f = \mathrm{id}_E$. Ce sont exactement les bijections de $E$ dans $E$, c'est-à-dire les permutations de $E$. Donc $\mathcal{U}_{E^E} = \mathfrak{S}_E$.

\begin{enumerate}
    \setcounter{enumi}{4}
    \item Soit $X$ un ensemble, et $\Delta : \mathcal{P}(X) \times \mathcal{P}(X) \to \mathcal{P}(X)$ la différence symétrique :
    \[
    (A, B) \longmapsto A \Delta B = (A \setminus B) \cup (B \setminus A) = (A \cup B) \setminus (A \cap B)
    \]
    $(\mathcal{P}(X), \Delta)$ est un groupe commutatif :
    \begin{itemize}
        \item \textbf{(A)} L'associativité a été prouvée dans le cours de théorie des ensembles.
        \item \textbf{(N)} $\emptyset$ est neutre pour $\Delta$ : $\forall A \in \mathcal{P}(X), \ A \Delta \emptyset = \emptyset \Delta A = A$.
        \item \textbf{(S)} $\forall A \in \mathcal{P}(X), \ A \Delta A = \emptyset$ (tout $A$ est symétrisable, égal à son symétrique).
    \end{itemize}
    
    \item Soit $(G, \cdot)$ un groupe et $I$ un ensemble non vide. L'ensemble $\mathcal{F}(I, G) = G^I$ peut être muni de la loi de composition interne :
    \[
    G^I \times G^I \longrightarrow G^I, \quad (f, g) \longmapsto f \cdot g
    \]
    où $f \cdot g : I \to G, \ x \mapsto f(x) \cdot g(x)$.
    
    Muni de cette loi, $G^I$ est un groupe :
    \begin{itemize}
        \item \textbf{(A)} Soient $f, g, h \in G^I$. Alors $\forall x \in I$ :
        \begin{align*}
            [(f \cdot g) \cdot h](x) &= (f \cdot g)(x) \cdot h(x) \\
                                     &= (f(x) \cdot g(x)) \cdot h(x) \\
                                     &= f(x) \cdot (g(x) \cdot h(x)) \quad (\text{associativité de } \cdot \text{ dans } G) \\
                                     &= [f \cdot (g \cdot h)](x)
        \end{align*}
        Donc $(f \cdot g) \cdot h = f \cdot (g \cdot h)$.
        \item \textbf{(N)} L'application constante $\chi : I \to G, \ x \mapsto e$ vérifie $\forall f \in G^I, \ f \cdot \chi = \chi \cdot f = f$.
        \item \textbf{(S)} Si $f : I \to G$, alors $f' : I \to G, \ x \mapsto f(x)^{-1}$ vérifie $f \cdot f' = f' \cdot f = \chi$.
    \end{itemize}
    \textit{Exemple particulier :} $(\mathcal{F}(I, \R), +)$ est un groupe dont le neutre est la fonction nulle.
    
    \item $(\U, \times)$ où $\U = \{ z \in \C \mid |z| = 1 \}$ est un groupe. Pour tout $n \in \N^*$, $(\U_n, \times)$, où $\U_n$ est l'ensemble des racines $n$-èmes de $1$, est un groupe.
    
    \item \textbf{Les groupes de la géométrie :}
    \begin{itemize}
        \item Groupe des isométries vectorielles du plan euclidien (ses éléments sont les rotations de centre $O$ et les réflexions vectorielles) ;
        \item Groupe des isométries du plan euclidien ;
        \item Même chose en remplaçant « plan » par « espace » ;
        \item Groupe des homographies de $\hat{\C} = \C \cup \{\infty\}$ (droite projective complexe, obtenue en ajoutant un point à l'infini à $\C$). Soient $(a, b, c, d) \in \C^4$ tel que $ad - bc \neq 0$. On note $h_{(a, b, c, d)} : \hat{\C} \to \hat{\C}$ :
        \begin{itemize}
            \item Si $c \neq 0$ :
            \[
            z \longmapsto \begin{cases}
            \dfrac{az+b}{cz+d} & \text{si } z \in \C \setminus \left\{-\dfrac{d}{c}\right\} \\[1.2em]
            \infty & \text{si } z = -\dfrac{d}{c} \\[0.6em]
            \dfrac{a}{c} & \text{si } z = \infty
            \end{cases}
            \]
            \item Si $c = 0$ (alors $a \neq 0$ et $d \neq 0$) :
            \[
            z \longmapsto \begin{cases}
            \dfrac{az+b}{d} & \text{si } z \in \C \\[1em]
            \infty & \text{si } z = \infty
            \end{cases}
            \]
        \end{itemize}
        \item Et les « petits » groupes : groupes des isométries d'un polygone régulier, d'un cube, d'un tétraèdre, d'un octaèdre, d'un dodécaèdre, d'un icosaèdre réguliers.
    \end{itemize}
\end{enumerate}
\end{exemples}

De l'étude des monoïdes, on sait que dans un groupe $(G, \cdot)$ :
\begin{itemize}
    \item L'élément neutre $e$ est unique ;
    \item Tout élément $x \in G$ a un symétrique unique et est régulier ; ce symétrique est noté $x'$ ou $x^{-1}$.
\end{itemize}

\begin{proposition}[Règles de calcul dans un groupe]
Soit $(G, \cdot)$ un groupe.
\begin{enumerate}
    \item Pour tous $x, y \in G$, $(xy)^{-1} = y^{-1} x^{-1}$.
    \item Pour tous $m, n \in \Z$ et tout $x \in G$, $x^{m+n} = x^m \cdot x^n$.
    \item Si $x$ et $y \in G$ commutent, alors pour tout $n \in \Z$, $(xy)^n = x^n y^n$.
\end{enumerate}
\end{proposition}

\begin{demonstration}
\begin{enumerate}
    \item Prouvé dans la partie I (théorème sur les monoïdes).
    
    \item Déjà établi pour $m, n \in \N$.
    \begin{itemize}
        \item \cas{Si $m \le 0$ et $n \le 0$ :} On a $-m, -n \in \N$, donc $x^{-m-n} = x^{-n} x^{-m}$. En passant aux inverses :
        \[
        \left[ x^{-(m+n)} \right]^{-1} = \left[ x^{-n} x^{-m} \right]^{-1} = (x^{-m})^{-1} (x^{-n})^{-1}
        \]
        soit $x^{m+n} = x^m \cdot x^n$.
        
        \item \cas{Si $m \le 0$ et $n \ge 0$ :}
        \begin{itemize}
            \item \textit{1er cas : $m + n \ge 0$.} Alors $x^{-m} \cdot x^{m+n} = x^{-m + m + n} = x^n$ (car $-m \ge 0$ et $m+n \ge 0$). Donc :
            \[
            x^{m+n} = (x^{-m})^{-1} x^n = x^m x^n
            \]
            \item \textit{2nd cas : $m + n \le 0$.} Alors $x^{m+n} \cdot x^{-n} = x^{m + n - n} = x^m$ (car $-(m+n) \ge 0$ et $-n \le 0$, en utilisant le cas précédent ou par passage à l'inverse sur $x^{-(m+n)} x^n = x^{-m}$). Donc :
            \[
            x^{m+n} = x^m (x^{-n})^{-1} = x^m x^n
            \]
        \end{itemize}
        
        \item \cas{Si $m \ge 0$ et $n \le 0$ :} Idem en échangeant les rôles de $m$ et $n$ (ou par commutativité de $m$ et $n$ dans $\Z$ et symétrie des deux cas précédents).
    \end{itemize}
    
    \item Déjà établi si $n \in \N$.
    Si $n \in \Z^-$, on a $-n = |n| \in \N$, d'où :
    \[
    (xy)^n = (xy)^{-|n|} = \left[ (xy)^{|n|} \right]^{-1} = \left[ x^{|n|} y^{|n|} \right]^{-1} = (y^{|n|})^{-1} (x^{|n|})^{-1} = y^n x^n
    \]
    Or $x$ et $y$ commutent, donc $x^n$ et $y^n$ commutent également (en effet, $xy = yx \implies y^{-1} x^{-1} = x^{-1} y^{-1}$, donc $x^{-1}$ et $y^{-1}$ commutent, et par récurrence leurs puissances positives $x^n = (x^{-1})^{|n|}$ et $y^n = (y^{-1})^{|n|}$ commutent). Ainsi :
    \[
    (xy)^n = x^n y^n
    \]
\end{enumerate}
\end{demonstration}

\begin{remarque}[Notation additive]
Si le groupe est commutatif, on peut noter additivement la loi de $G$ (c'est-à-dire par $+$). Le neutre se note $0$ et le symétrique de $x$ se note $-x$.

Dans ce cas, $x^n$ se note $n \cdot x$ (attention : $\cdot$ n'est pas une loi de multiplication interne, mais une loi externe $\Z \times G \to G$), défini par :
\[
n \cdot x = \begin{cases}
\underbrace{x + x + \dots + x}_{n \text{ fois}} & \text{si } n \ge 1 \\[1.2em]
0 & \text{si } n = 0 \\[0.4em]
(-x) + (-x) + \dots + (-x) \quad (|n| \text{ fois}) & \text{si } n < 0
\end{cases}
\]
Les propriétés 1), 2), 3) s'écrivent alors :
\begin{enumerate}
    \item $\forall x, y \in G, \quad -(x + y) = -x - y$ \quad (c'est-à-dire $-x + (-y)$)
    \item $\forall m, n \in \Z, \ \forall x \in G, \quad (m + n) \cdot x = m \cdot x + n \cdot x$
    \item $\forall x, y \in G, \ \forall n \in \Z, \quad n \cdot (x + y) = n \cdot x + n \cdot y$
\end{enumerate}
\end{remarque}

\begin{remarque}
Si $x \in G$ et $y \in G$ vérifient $x \cdot y = e$, alors $y = x^{-1}$.
En effet, $x$ étant inversible :
\[
x \cdot y = e \implies x^{-1} \cdot (x \cdot y) = x^{-1} \cdot e \implies (x^{-1} \cdot x) \cdot y = x^{-1} \implies e \cdot y = x^{-1} \implies y = x^{-1}
\]
De même, $y \cdot x = e \implies y = x^{-1}$.
\end{remarque}

\newpage
\subsection{2.2. Translations dans un groupe}

\begin{proposition}
Soit $(G, \cdot)$ un groupe. Les translations à droite et à gauche sont des permutations de $G$.
\end{proposition}

\begin{demonstration}
Soit $a \in G$ et $\delta_a : G \to G, \ x \mapsto x \cdot a$.
Soit $y \in G$. Il s'agit de montrer que :
\[
\exists! x \in G, \quad x \cdot a = y
\]
\begin{itemize}
    \item \cas{Analyse :} Si $x \in G$ vérifie $x \cdot a = y$, alors $(x \cdot a) \cdot a^{-1} = y \cdot a^{-1}$, soit $x \cdot (a \cdot a^{-1}) = y \cdot a^{-1}$, et donc $x = y \cdot a^{-1}$.
    \item \cas{Synthèse :} Réciproquement, en posant $x = y \cdot a^{-1}$, on a :
    \[
    (y \cdot a^{-1}) \cdot a = y \cdot (a^{-1} \cdot a) = y \cdot e = y
    \]
\end{itemize}
Donc $\delta_a : G \to G$ est une bijection, et $(\delta_a)^{-1} : G \to G, \ y \mapsto y \cdot a^{-1}$, c'est-à-dire :
\[
(\delta_a)^{-1} = \delta_{a^{-1}}
\]
De même, $\gamma_a : G \to G, \ x \mapsto a \cdot x$ est une bijection et :
\[
(\gamma_a)^{-1} = \gamma_{a^{-1}}
\]
\end{demonstration}

\textbf{Table de Cayley d'un groupe :}
On peut définir un groupe fini par la table de sa loi $\cdot$ sur $G$.

\begin{exemples}[Table de Cayley d'un groupe d'ordre 3]
\[
\begin{array}{c|ccc}
\cdot & a & b & c \\
\hline
a & a & b & c \\
b & b & c & a \\
c & c & a & b
\end{array}
\]
C'est la table d'une loi de groupe dont le neutre est $a$.
\end{exemples}

En notant $G = \{a_1, a_2, \dots, a_n\}$, la ligne n°$i$ de la table de Cayley de $G$ représente les éléments $\gamma_{a_i}(a_1), \gamma_{a_i}(a_2), \dots, \gamma_{a_i}(a_n)$ dans cet ordre, et la colonne n°$j$ de la table représente les éléments $\delta_{a_j}(a_1), \delta_{a_j}(a_2), \dots, \delta_{a_j}(a_n)$ dans cet ordre (de haut en bas).

Comme les translations (à droite et à gauche) sont des permutations de $G$, une condition nécessaire pour qu'une table de loi soit celle d'un groupe est que toute ligne (resp. colonne) de la table soit une permutation de l'ensemble des éléments de $G$.

Cette condition n'est cependant \mot{pas suffisante} :

\begin{exemples}[Contre-exemple]
\[
\begin{array}{c|ccc}
\cdot & a & b & c \\
\hline
a & a & b & c \\
b & c & a & b \\
c & b & c & a
\end{array}
\]
Chaque ligne et chaque colonne est bien une permutation de $\{a, b, c\}$, mais cette loi ne possède pas de neutre (donc ce n'est pas un groupe).
\end{exemples}

\newpage
\subsection{2.3. Groupe produit}

\begin{proposition}
Soient $(G_1, *)$ et $(G_2, \otimes)$ deux groupes. L'ensemble $G_1 \times G_2$ peut être muni de la loi de composition interne $T$ définie par :
\[
(x_1, x_2) \, T \, (y_1, y_2) = (x_1 * y_1, \, x_2 \otimes y_2)
\]
où $(x_1, x_2), (y_1, y_2) \in G_1 \times G_2$.
Alors $(G_1 \times G_2, T)$ est un groupe, appelé \mot{groupe produit} de $(G_1, *)$ et $(G_2, \otimes)$.
\end{proposition}

\begin{demonstration}
D'abord, $T$ est bien une loi de composition interne sur $G_1 \times G_2$. Notons $e_1$ (resp. $e_2$) le neutre de $G_1$ (resp. $G_2$).
\begin{itemize}
    \item \textbf{(N)} On vérifie aisément que $(e_1, e_2)$ est neutre pour $T$.
    \item \textbf{(A) Associativité :} Soient $(x_1, x_2), (y_1, y_2)$ et $(z_1, z_2)$ trois éléments de $G_1 \times G_2$. Alors :
    \begin{align*}
        \big((x_1, x_2) \, T \, (y_1, y_2)\big) \, T \, (z_1, z_2) 
        &= (x_1 * y_1, \, x_2 \otimes y_2) \, T \, (z_1, z_2) \\
        &= \big((x_1 * y_1) * z_1, \, (x_2 \otimes y_2) \otimes z_2\big) \\
        &= \big(x_1 * (y_1 * z_1), \, x_2 \otimes (y_2 \otimes z_2)\big) \\
        &= (x_1, x_2) \, T \, \big((y_1, y_2) \, T \, (z_1, z_2)\big)
    \end{align*}
    \item \textbf{(S)} Soit $(x_1, x_2) \in G_1 \times G_2$. Notons $x_1'$ le symétrique de $x_1$ dans $G_1$ et $x_2'$ le symétrique de $x_2$ dans $G_2$. Alors :
    \[
    (x_1, x_2) \, T \, (x_1', x_2') = (x_1 * x_1', \, x_2 \otimes x_2') = (e_1, e_2)
    \]
    et
    \[
    (x_1', x_2') \, T \, (x_1, x_2) = (x_1' * x_1, \, x_2' \otimes x_2) = (e_1, e_2)
    \]
\end{itemize}
Donc $(G_1 \times G_2, T)$ est un groupe.
\end{demonstration}

\textbf{Généralisation :} Si $(G_k, *_k)_{k \in \intEnt{1}{n}}$ est une famille de $n$ groupes, le produit cartésien $G_1 \times G_2 \times \dots \times G_n = \displaystyle\prod_{k=1}^n G_k$ est un groupe pour la loi produit :
\[
(x_1, x_2, \dots, x_n) \cdot (y_1, y_2, \dots, y_n) = (x_1 *_1 y_1, \dots, x_n *_n y_n)
\]

\begin{exemples}[Groupe de Klein]
$\U_2 = \{1, -1\}$ est un groupe à $2$ éléments (de neutre $1$).
Notons que si $G = \{e, a\}$ est un groupe à $2$ éléments de neutre $e$, sa table de Cayley est :
\[
\begin{array}{c|cc}
\cdot & e & a \\
\hline
e & e & a \\
a & a & e
\end{array}
\qquad \text{avec } a \cdot a = e \text{ (c'est-à-dire } a^2 = e\text{)}.
\]
Réciproquement, cette table définit bien une loi de groupe.

Formons la table de Cayley du groupe produit $(\U_2 \times \U_2, \times)$ (\mot{groupe de Klein}) :
\[
\begin{array}{c|cccc}
\times & (1, 1) & (1, -1) & (-1, 1) & (-1, -1) \\
\hline
(1, 1) & (1, 1) & (1, -1) & (-1, 1) & (-1, -1) \\
(1, -1) & (1, -1) & (1, 1) & (-1, -1) & (-1, 1) \\
(-1, 1) & (-1, 1) & (-1, -1) & (1, 1) & (1, -1) \\
(-1, -1) & (-1, -1) & (-1, 1) & (1, -1) & (1, 1)
\end{array}
\]
On vérifie que $\forall x \in \U_2 \times \U_2, \quad x^2 = (1, 1)$ (neutre de $\U_2 \times \U_2$).

Au contraire, le groupe $(\U_4, \times)$ ne vérifie pas cette propriété :
\[
\U_4 = \{1, i, -1, -i\} \quad (\text{neutre } 1)
\]
\[
i^2 = -1 \neq 1 \quad \text{et} \quad (-i)^2 = -1 \neq 1
\]
Les tables de Cayley de ces deux groupes $(\U_2 \times \U_2, \times)$ et $(\U_4, \times)$ sont fondamentalement différentes : on dira qu'ils ne sont pas \mot{isomorphes}.
\end{exemples}

\newpage
\section{III -- Sous-groupes d'un groupe}

\subsection{3.1. Définitions et caractérisation}

\begin{definition}
Soit $(G, \cdot)$ un groupe et $H \subset G$. On dit que $H$ est un \mot{sous-groupe} de $G$ si $H$ est stable par la loi $\cdot$, et $(H, \cdot)$ est un groupe pour la loi induite.
\end{definition}

\begin{lemme}
Soit $H$ un sous-groupe de $(G, \cdot)$. Alors $e_H = e_G$, et les inverses d'éléments $h \in H$ dans $H$ et $G$ coïncident.
\end{lemme}

\begin{demonstration}
On note $e_G$ le neutre de $(G, \cdot)$ et $e_H$ le neutre de $(H, \cdot)$. Alors :
\[
e_H \cdot e_H = e_H \quad (\text{car } e_H \text{ est neutre dans } (H, \cdot))
\]
\[
e_H = e_G \cdot e_H \quad (\text{car } e_G \text{ est neutre dans } (G, \cdot))
\]
c'est-à-dire $e_H \cdot e_H = e_G \cdot e_H$. Or dans le groupe $(G, \cdot)$, on peut simplifier à droite par $e_H$ (tout élément d'un groupe étant régulier), d'où il vient :
\[
\cadre{e_H = e_G}
\]

Soit $h \in H$, $h'$ son inverse dans $H$ et $h^{-1}$ son inverse dans $G$. Alors :
\[
h \cdot h' = e_H = e_G = h \cdot h^{-1}
\]
Donc, en simplifiant à gauche par $h$ dans le groupe $(G, \cdot)$, il vient :
\[
\cadre{h' = h^{-1}}
\]
\end{demonstration}

\begin{proposition}[Caractérisation des sous-groupes]
Soit $(G, \cdot)$ un groupe et $H \subset G$. Alors $H$ est un sous-groupe de $(G, \cdot)$ si et seulement si l'une des deux conditions équivalentes suivantes est vérifiée :
\begin{enumerate}
    \item \begin{itemize}
        \item $e_G \in H$
        \item $\forall (x, y) \in H^2, \quad x \cdot y \in H$
        \item $\forall x \in H, \quad x^{-1} \in H$ \quad (où $x^{-1}$ est l'inverse dans $G$)
    \end{itemize}
    \item \begin{itemize}
        \item $H \neq \emptyset$
        \item $\forall (x, y) \in H^2, \quad x^{-1} \cdot y \in H$ \quad (ou $\forall (x, y) \in H^2, \quad x \cdot y^{-1} \in H$)
    \end{itemize}
\end{enumerate}
\end{proposition}

\begin{demonstration}
\cas{Équivalence entre « $H$ est un sous-groupe de $(G, \cdot)$ » et 1) :}
\begin{itemize}
    \item \textbf{Sens direct ($\implies$) :} Supposons que $H$ soit un sous-groupe de $(G, \cdot)$.
    D'après le lemme, $e_G = e_H \in H$.
    Par définition, $H$ est stable par $\cdot$, donc $\forall (x, y) \in H^2, \ x \cdot y \in H$.
    Le lemme dit aussi que si $x \in H$, alors son inverse dans $G$ coïncide avec son inverse dans $H$, donc $x^{-1} \in H$.
    Ainsi, la condition 1) est vérifiée.
    
    \item \textbf{Réciproque ($\Longleftarrow$) :} Supposons 1).
    Alors $H$ est stable par la loi $\cdot$ de $G$. Vérifions que $H$, muni de la loi induite, est bien un groupe :
    \begin{itemize}
        \item \textbf{(A)} L'associativité de la loi induite sur $H$ découle de l'associativité de la loi de $G$.
        \item \textbf{(N)} $e_G \in H$ et $\forall x \in H, \ x \cdot e_G = x = e_G \cdot x$, donc $H$ possède un neutre, à savoir $e_G$.
        \item \textbf{(S)} Soit $x \in H$. Alors $x^{-1} \in H$ et $x \cdot x^{-1} = x^{-1} \cdot x = e_G$, où $e_G$ est le neutre de $G$, mais aussi celui de $H$. Donc tout élément de $H$ est symétrisable pour la loi induite.
    \end{itemize}
    Donc $(H, \cdot)$ est un groupe, c'est-à-dire un sous-groupe de $(G, \cdot)$.
\end{itemize}

\vspace{8pt}
\cas{Équivalence entre 1) et 2) :}
\begin{itemize}
    \item \textbf{1) $\implies$ 2) :} Supposons 1). Comme $e_G \in H$, on a $H \neq \emptyset$.
    De plus, soient $(x, y) \in H^2$. Par 1), $x^{-1} \in H$, et par stabilité de $H$, $x^{-1} \cdot y \in H$. D'où 2).
    
    \item \textbf{2) $\implies$ 1) :} Supposons 2).
    Comme $H \neq \emptyset$, on peut choisir $x_0 \in H$. Alors $x_0^{-1} \cdot x_0 \in H$, c'est-à-dire $e_G \in H$.
    
    Dans la propriété $\forall (x, y) \in H^2, \ x^{-1} \cdot y \in H$, on peut remplacer $y$ par $e_G \in H$, ce qui donne :
    \[
    \forall x \in H, \quad x^{-1} \cdot e_G = x^{-1} \in H
    \]
    Alors, pour tout $(x, y) \in H^2$, on a $(x^{-1}, y) \in H^2$, et donc $(x^{-1})^{-1} \cdot y \in H$, c'est-à-dire $x \cdot y \in H$.
    Donc 2) $\implies$ 1).
\end{itemize}
\end{demonstration}

\begin{remarque}[Caractérisation d'un intérêt pratique important]
Dans de nombreux cas, pour montrer qu'un magma $(G, \cdot)$ est un groupe, on montre que c'est un sous-groupe d'un groupe bien connu.
\end{remarque}

\begin{exemples}
\begin{enumerate}
    \item $\{e\}$ et $G$ sont des sous-groupes de $(G, \cdot)$ (dits \mot{sous-groupes triviaux}).
    
    \item $\U = \{ z \in \C \mid |z| = 1 \}$ est un sous-groupe de $(\C^*, \times)$ :
    \begin{itemize}
        \item D'abord, $\U \subset \C^*$ ;
        \item $1 \in \U$ (car $|1| = 1$) ;
        \item $\forall z, z' \in \U, \ |zz'| = |z||z'| = 1 \times 1 = 1$, donc $zz' \in \U$ ;
        \item $\forall z \in \U, \ z^{-1} = \bar{z}$, donc $|z^{-1}| = |\bar{z}| = |z| = 1$, donc $z^{-1} \in \U$.
    \end{itemize}
    Donc $\U$ est un sous-groupe de $(\C^*, \times)$.
    
    \item Pour tout $n \in \N^*$, $\U_n = \{ z \in \C \mid z^n = 1 \}$ est un sous-groupe de $(\C^*, \times)$ (et de $(\U, \times)$) :
    \begin{itemize}
        \item D'abord, $\U_n \subset \C^*$ (en fait $\U_n \subset \U \subset \C^*$) ;
        \item $1 \in \U_n$ (car $1^n = 1$) ;
        \item Si $z, z' \in \U_n$, alors $(zz')^n = z^n z'^n = 1 \times 1 = 1$, donc $zz' \in \U_n$ ;
        \item Si $z \in \U_n$, $(z^{-1})^n = \dfrac{1}{z^n} = \dfrac{1}{1} = 1$, donc $z^{-1} \in \U_n$.
    \end{itemize}
    Donc $\U_n$ est un sous-groupe de $(\C^*, \times)$ (et de $(\U, \times)$).
    
    \item $\U_\infty = \{ z \in \C \mid \exists n \in \N^*, \ z^n = 1 \} = \displaystyle\bigcup_{n \ge 1} \U_n$ (groupe des racines de l'unité) est un sous-groupe de $(\C^*, \times)$ (et de $(\U, \times)$) :
    \begin{itemize}
        \item $1 \in \U_\infty$ (car $1^1 = 1$) ;
        \item Si $z, z' \in \U_\infty$, alors il existe $n$ et $n' \ge 1$ tels que $z^n = 1$ et $z'^{n'} = 1$. Alors (comme $nn' \ge 1$) :
        \[
        (z^{-1} z')^{nn'} = (z^{-1})^{nn'} (z')^{nn'} = \left[ (z^{-1})^n \right]^{n'} \left[ (z')^{n'} \right]^n = \left( \frac{1}{z^n} \right)^{n'} \times 1^n = 1^{n'} \times 1^n = 1
        \]
        Donc $z^{-1} z' \in \U_\infty$.
    \end{itemize}
\end{enumerate}
\end{exemples}

\begin{proposition}[Intersection de sous-groupes]
Soit $(H_i)_{i \in I}$ une famille non vide ($I \neq \emptyset$) de sous-groupes de $(G, \cdot)$. Alors $\displaystyle\bigcap_{i \in I} H_i$ est un sous-groupe de $(G, \cdot)$.
\end{proposition}

\begin{demonstration}
D'abord, $\displaystyle\bigcap_{i \in I} H_i \subset G$.
\begin{enumerate}
    \item $\forall i \in I, \ e \in H_i$ car $H_i$ est un sous-groupe de $(G, \cdot)$. Donc $e \in \displaystyle\bigcap_{i \in I} H_i$ (en particulier, $\displaystyle\bigcap_{i \in I} H_i \neq \emptyset$).
    \item Soient $x, y \in \displaystyle\bigcap_{i \in I} H_i$. Alors $\forall i \in I, \ x \in H_i$ et $y \in H_i$, donc $x^{-1} \cdot y \in H_i$.
    Ainsi, $x^{-1} \cdot y \in \displaystyle\bigcap_{i \in I} H_i$.
\end{enumerate}
De 1) et 2), il résulte que $\displaystyle\bigcap_{i \in I} H_i$ est un sous-groupe de $(G, \cdot)$.
\end{demonstration}

\begin{remarque}
\textbf{Attention :} La réunion de deux sous-groupes n'est pas toujours un sous-groupe.
\end{remarque}

\begin{proposition}[Réunion de deux sous-groupes -- Question de cours]
Soient $H$ et $K$ deux sous-groupes de $(G, \cdot)$. Alors $H \cup K$ est un sous-groupe de $(G, \cdot)$ si et seulement si $H \subset K$ ou $K \subset H$.
\end{proposition}

\begin{demonstration}
\begin{itemize}
    \item \textbf{Sens direct ($\implies$) :} Supposons que $H \cup K$ soit un sous-groupe de $(G, \cdot)$. Montrons que $H \subset K$ ou $K \subset H$, ce qui revient à montrer que $\neg(H \subset K) \implies K \subset H$.
    
    Supposons $H \not\subset K$ ; il existe donc $h_0 \in H \setminus K$. Montrons que $K \subset H$.
    
    Soit $k \in K$. Comme $h_0$ et $k$ appartiennent à $H \cup K$ qui est un sous-groupe de $(G, \cdot)$, on a $h_0 \cdot k \in H \cup K$.
    
    Or $h_0 \cdot k \notin K$, car sinon on aurait :
    \[
    h_0 = \underbrace{(h_0 \cdot k)}_{\in K} \cdot \underbrace{k^{-1}}_{\in K} \in K
    \]
    ce qui est contradictoire avec $h_0 \notin K$.
    
    Donc $h_0 \cdot k \in H$, d'où :
    \[
    k = \underbrace{h_0^{-1}}_{\in H} \cdot \underbrace{(h_0 \cdot k)}_{\in H} \in H
    \]
    Ainsi $K \subset H$.
    
    \item \textbf{Réciproque ($\Longleftarrow$) :} Si $H \subset K$ ou $K \subset H$, alors $H \cup K = K$ ou $H \cup K = H$, qui sont par hypothèse déjà des sous-groupes de $(G, \cdot)$.
\end{itemize}
\end{demonstration}

\begin{remarque}[Exercices sur les réunions de sous-groupes]
\begin{enumerate}
    \item Soit $(H_n)_{n \in \N}$ une suite croissante de sous-groupes de $(G, \cdot)$. Alors $\displaystyle\bigcup_{n \in \N} H_n$ est un sous-groupe de $(G, \cdot)$.
    \item La suite $(\U_n)_{n \in \N^*}$ n'est pas croissante (pour l'inclusion), mais $\forall m, n \in \N^*, \ \exists p \in \N^*$ tel que $\U_m \cup \U_n \subset \U_p$ (par exemple $p = mn$). Cette propriété permet également de montrer que $\U_\infty = \displaystyle\bigcup_{n \ge 1} \U_n$ est un sous-groupe de $(\U, \times)$.
\end{enumerate}
\end{remarque}

\newpage
\subsection{3.2. Sous-groupe engendré par une partie}

\begin{proposition}
Soient $(G, \cdot)$ un groupe et $S \subset G$. Il existe un plus petit sous-groupe (pour l'inclusion) contenant $S$. C'est l'intersection de l'ensemble des sous-groupes de $(G, \cdot)$ contenant $S$. On l'appelle le \mot{sous-groupe engendré par $S$} et on le note $\langle S \rangle$.
\end{proposition}

\begin{demonstration}
Notons $\mathscr{S} = \{ H \mid H \text{ sous-groupe de } G \text{ et } S \subset H \}$.
$\mathscr{S} \neq \emptyset$ car $G \in \mathscr{S}$. On peut donc former $\displaystyle\bigcap_{H \in \mathscr{S}} H$, qui est un sous-groupe de $(G, \cdot)$ d'après le paragraphe 3.1.

De plus, pour tout $H \in \mathscr{S}$, $S \subset H$, donc $S \subset \displaystyle\bigcap_{H \in \mathscr{S}} H$.

Enfin, si $H_0 \in \mathscr{S}$, alors $\displaystyle\bigcap_{H \in \mathscr{S}} H \subset H_0$.

Donc $\displaystyle\bigcap_{H \in \mathscr{S}} H$ est le plus petit sous-groupe de $G$ contenant $S$.
\end{demonstration}

\begin{remarque}
Si $S = \emptyset$, alors $\langle \emptyset \rangle = \{e\}$ (car $\emptyset \subset \{e\}$ est vraie).
\end{remarque}

\begin{proposition}[Sous-groupe engendré par un élément]
Soit $(G, \cdot)$ un groupe et $x \in G$. Le sous-groupe engendré par $\{x\}$ (noté $\langle x \rangle$) est :
\[
\cadre{\langle x \rangle = \{ x^n \mid n \in \Z \}}
\]
\end{proposition}

\begin{demonstration}
Par définition, $\langle x \rangle$ contient $e$ et $x$. Comme $\langle x \rangle$ est un sous-groupe, on a aussi $x^{-1} \in \langle x \rangle$.
Par stabilité du sous-groupe $\langle x \rangle$ par la loi $\cdot$, on a :
\[
\forall n \ge 1, \quad x^n \in \langle x \rangle \quad \text{et} \quad (x^{-1})^n = x^{-n} \in \langle x \rangle
\]
Donc $\forall n \in \Z^*, \ x^n \in \langle x \rangle$. Comme $x^0 = e \in \langle x \rangle$, on a bien :
\[
\{ x^n \mid n \in \Z \} \subset \langle x \rangle
\]

Il suffit de vérifier que $H = \{ x^n \mid n \in \Z \}$ est un sous-groupe de $(G, \cdot)$ pour conclure par minimalité de $\langle x \rangle$ :
\begin{itemize}
    \item $x^0 = e \in H$ ;
    \item Soient $m, n \in \Z$. Alors $x^m \cdot x^n = x^{m+n} \in H$ ;
    \item Si $m \in \Z$, alors $(x^m)^{-1} = x^{-m} \in H$.
\end{itemize}
Donc $\{ x^n \mid n \in \Z \}$ est un sous-groupe de $(G, \cdot)$ contenant $x$, d'où $\langle x \rangle = \{ x^n \mid n \in \Z \}$.
\end{demonstration}

\begin{definition}[Groupe monogène]
Un groupe $(G, \cdot)$ est dit \mot{monogène} s'il est engendré par un seul élément, c'est-à-dire s'il existe $a \in G$ tel que :
\[
G = \langle a \rangle = \{ a^n \mid n \in \Z \}
\]
Un tel élément $a$ s'appelle un \mot{générateur} du groupe $(G, \cdot)$ (on dit aussi que $a$ engendre $(G, \cdot)$).
\end{definition}

\begin{exemples}
\begin{enumerate}
    \item $(\U_n, \times)$ est un groupe monogène engendré par $\omega = e^{\frac{2i\pi}{n}}$ (on montre que les générateurs de $\U_n$ sont exactement les $\omega^k$ où $1 \le k \le n$ et $\mathrm{pgcd}(k, n) = 1$).
    
    \item $(\Z, +)$ est un groupe monogène. Attention au passage en notation additive : si $a \in (G, +)$, alors $\langle a \rangle = \{ n \cdot a \mid n \in \Z \}$. En effet :
    \[
    \Z = \langle 1 \rangle = \langle -1 \rangle
    \]
    
    \item $(\Z^2, +)$ n'est pas monogène. En effet, si $(m_0, n_0) \in \Z^2 \setminus \{(0, 0)\}$, le sous-groupe engendré par $(m_0, n_0)$ est :
    \[
    \langle (m_0, n_0) \rangle = \{ k \cdot (m_0, n_0) \mid k \in \Z \} = \{ (k m_0, \, k n_0) \mid k \in \Z \}
    \]
    Un couple $(x, y) \in \Z^2$ tel que $n_0 x \neq m_0 y$ n'appartient pas à ce sous-groupe (par exemple, si $m_0 \neq 0$, le couple $(0, 1)$ n'y appartient pas ; si $n_0 \neq 0$, $(1, 0)$ n'y appartient pas).
    
    \item \textit{Généralisation (Exercice) :} Soit $(G, +)$ un groupe commutatif noté additivement et $S = \{a_1, \dots, a_n\} \subset G$. Alors le sous-groupe engendré par $S$ est :
    \[
    \langle a_1, \dots, a_n \rangle = \{ x \in G \mid \exists (k_1, \dots, k_n) \in \Z^n, \ x = k_1 \cdot a_1 + \dots + k_n \cdot a_n \}
    \]
\end{enumerate}
\end{exemples}

\newpage
\subsection{3.3. Les sous-groupes de $(\Z, +)$}

\textbf{a) Division euclidienne dans $(\Z, +)$}

\begin{theoreme}[Division euclidienne dans $\Z$]
Pour tout $(a, b) \in \Z \times \N^*$, il existe un unique couple $(q, r) \in \Z \times \N$ tel que :
\[
\begin{cases}
a = qb + r \\
0 \le r < b
\end{cases}
\]
\end{theoreme}

\begin{demonstration}
\cas{Existence :} Déjà établie lorsque $a \in \N$.

Supposons $a \in \Z^-$. Alors $a + b|a| \ge a + |a| \ge 0$ (puisque $b \ge 1$), donc $a + b|a| \in \N$.

D'après le cas positif, il existe $(q_0, r) \in \N^2$ tel que :
\[
a + b|a| = q_0 b + r \quad \text{et} \quad 0 \le r < b
\]
En posant $q = q_0 - |a|$, on a $q \in \Z$ et :
\[
a = qb + r \quad \text{avec} \quad 0 \le r < b
\]

\cas{Unicité :} Supposons qu'il existe $(q, r)$ et $(q', r') \in \Z \times \intEnt{0}{b-1}$ tels que $a = bq + r$ et $a = bq' + r'$.
Alors :
\[
b(q - q') = r' - r \implies b|q - q'| = |r' - r| < b
\]
Donc $|q' - q| < 1$. Comme $|q' - q| \in \N$, on a $|q' - q| = 0$, d'où $q' = q$, et par suite $r = r'$.
\end{demonstration}

\vspace{8pt}
\textbf{b) Sous-groupes de $(\Z, +)$}

\begin{theoreme}
Les sous-groupes de $(\Z, +)$ sont les parties de $\Z$ de la forme $m\Z$ où $m \in \N$.
\end{theoreme}

\begin{remarque}
Ce théorème affirme que tous les sous-groupes de $(\Z, +)$ sont monogènes. Si $m \in \N$, $m\Z = \{ k \cdot m \mid k \in \Z \} = \langle m \rangle$ (et on a aussi $\langle -m \rangle = \langle m \rangle$).
\end{remarque}

\begin{demonstration}
\begin{itemize}
    \item \textbf{Sens direct :} D'abord, si $m \in \N$, $m\Z$ est un sous-groupe de $(\Z, +)$ puisque c'est le sous-groupe engendré par $m$.
    
    \item \textbf{Réciproque :} Soit $H$ un sous-groupe de $(\Z, +)$.
    \begin{itemize}
        \item Si $H = \{0\}$, alors $H = \langle 0 \rangle = 0\Z$.
        \item Supposons $H \neq \{0\}$. Alors il existe $x_0 \in \Z^*$ tel que $x_0 \in H$. Comme $H$ est un sous-groupe de $(\Z, +)$, on a aussi $-x_0 \in H$, donc $|x_0| \in H \cap \N^*$.
        
        Ainsi, $H \cap \N^*$ est une partie non vide de $\N^*$ ; à ce titre, elle possède un plus petit élément, noté $m = \min(H \cap \N^*)$.
        
        Montrons que $H = m\Z$ par double inclusion :
        \begin{itemize}
            \item \textit{Inclusion $m\Z \subset H$ :} $m\Z = \langle m \rangle$ est le sous-groupe engendré par $m$, et $m \in H$ qui est un sous-groupe. Donc $\langle m \rangle \subset H$, c'est-à-dire $m\Z \subset H$.
            \item \textit{Inclusion $H \subset m\Z$ :} Soit $x \in H$. Comme $m \neq 0$, on peut diviser euclidiennement $x$ par $m$ : il existe $q \in \Z$ et $r \in \intEnt{0}{m-1}$ tels que :
            \[
            x = q \cdot m + r
            \]
            Mais alors $r = x - qm \in H$ (car $x \in H$ et $qm \in H$). De plus, $r \ge 0$, donc $r \in H \cap \N$.
            Or $r < m = \min(H \cap \N^*)$, donc $r \in (H \cap \N) \setminus (H \cap \N^*)$, ce qui impose $r = 0$.
            Ainsi $x = qm \in m\Z$, ce qui prouve $H \subset m\Z$.
        \end{itemize}
    \end{itemize}
\end{itemize}
Par double inclusion, $H = m\Z$.
\end{demonstration}

\begin{remarque}
Dans $\Z$, la notation $q \cdot m = q \times m$ désigne à la fois la loi externe $\Z \times \Z \to \Z$ et la loi interne de multiplication usuelle.
\end{remarque}

\newpage
\subsection{3.4. Classes à droite / Classes à gauche modulo un sous-groupe -- Théorème de Lagrange}

\begin{definition}[Classes à gauche et à droite modulo un sous-groupe]
Soit $(G, \cdot)$ un groupe et $H \subset G$ un sous-groupe de $(G, \cdot)$.
\begin{enumerate}
    \item La relation binaire $\mathcal{R}_g$ sur $G$ définie par :
    \[
    \forall (x, y) \in G^2, \quad x \, \mathcal{R}_g \, y \iff x^{-1} y \in H
    \]
    est une relation d'équivalence sur $G$, dont les classes d'équivalence sont les parties de $G$ de la forme $xH$ où $x \in G$. Elles sont appelées \mot{classes à gauche} de $G$ modulo $H$.
    
    L'ensemble quotient est noté $G / H$ :
    \[
    G / H = \{ xH \mid x \in G \}
    \]
    
    \item La relation binaire $\mathcal{R}_d$ sur $G$ définie par :
    \[
    \forall (x, y) \in G^2, \quad x \, \mathcal{R}_d \, y \iff x y^{-1} \in H
    \]
    est une relation d'équivalence sur $G$, dont les classes d'équivalence sont les parties de $G$ de la forme $Hx$ où $x \in G$. Elles sont appelées \mot{classes à droite} de $G$ modulo $H$.
    
    L'ensemble quotient est noté $H \backslash G$ :
    \[
    H \backslash G = \{ Hx \mid x \in G \}
    \]
\end{enumerate}
\end{definition}

\begin{lemme}
Soit $H$ un sous-groupe de $(G, \cdot)$ et $(x, y) \in G^2$. Alors :
\begin{enumerate}
    \item $x^{-1} y \in H \iff xH = yH$
    \item $x y^{-1} \in H \iff Hx = Hy$
\end{enumerate}
\end{lemme}

\begin{demonstration}
\begin{enumerate}
    \item \begin{itemize}
        \item \textbf{Sens direct ($\implies$) :} Supposons $x^{-1} y \in H$, donc $x^{-1} y = h$ où $h \in H$, d'où $y = xh$.
        Alors :
        \[
        yH = (xh)H = x(hH)
        \]
        c'est-à-dire $yH = \gamma_y(H) = \gamma_{xh}(H) = (\gamma_x \circ \gamma_h)(H) = \gamma_x(\gamma_h(H))$.
        Or, $h$ étant un élément du sous-groupe $H$, la translation $\gamma_h$ induit une permutation de $H$, donc $\gamma_h(H) = H$ (soit $hH = H$).
        Finalement, $yH = \gamma_x(H) = xH$.
        
        \item \textbf{Réciproque ($\Longleftarrow$) :} Supposons $xH = yH$. En particulier, $y = y \cdot e \in yH = xH$, donc il existe $h \in H$ tel que $y = xh$, d'où $x^{-1} y = h \in H$.
    \end{itemize}
    
    \item \textbf{Preuve analogue} (en utilisant la translation à droite) :
    Si $xy^{-1} = h \in H$, alors $x = hy$, et en utilisant $\delta_x = \delta_{hy} = \delta_y \circ \delta_h$ :
    \[
    Hx = \delta_x(H) = \delta_{hy}(H) = (\delta_y \circ \delta_h)(H) = \delta_y(\delta_h(H)) = \delta_y(H) = Hy
    \]
    et réciproquement si $Hx = Hy$, alors $x = e \cdot x \in Hx = Hy$, donc $x = hy$ avec $h \in H$, soit $xy^{-1} = h \in H$.
\end{enumerate}
\end{demonstration}

\begin{demonstration}[Preuve de la proposition sur $\mathcal{R}_g$ et $\mathcal{R}_d$]
D'après le lemme, on a pour tout $(x, y) \in G^2$ :
\[
x \, \mathcal{R}_g \, y \iff x^{-1} y \in H \iff xH = yH
\]
Il en résulte aisément que $\mathcal{R}_g$ est une relation d'équivalence sur $G$ :
\begin{itemize}
    \item \textbf{Réflexivité :} $\forall x \in G, \ x \, \mathcal{R}_g \, x$ car $xH = xH$.
    \item \textbf{Symétrie :} $\forall x, y \in G, \ x \, \mathcal{R}_g \, y \iff xH = yH \iff yH = xH \iff y \, \mathcal{R}_g \, x$.
    \item \textbf{Transitivité :} $\forall x, y, z \in G, \ (x \, \mathcal{R}_g \, y \text{ et } y \, \mathcal{R}_g \, z) \implies (xH = yH \text{ et } yH = zH) \implies xH = zH \implies x \, \mathcal{R}_g \, z$.
\end{itemize}
Soit $x \in G$. La classe d'équivalence de $x$ modulo $\mathcal{R}_g$ est :
\begin{align*}
    \mathrm{Cl}_{\mathcal{R}_g}(x) &= \{ y \in G \mid x \, \mathcal{R}_g \, y \} \\
                                   &= \{ y \in G \mid x^{-1} y \in H \} \\
                                   &= \{ y \in G \mid \exists h \in H, \ x^{-1} y = h \} \\
                                   &= \{ y \in G \mid \exists h \in H, \ y = xh \} \\
                                   &= \gamma_x(H) = xH
\end{align*}
On prouve de même que $\mathcal{R}_d$ est une relation d'équivalence sur $G$ et que pour tout $x \in G$, $\mathrm{Cl}_{\mathcal{R}_d}(x) = Hx$.
\end{demonstration}

\begin{proposition}[Équipotence des ensembles quotients $G / H$ et $H \backslash G$]
Soit $H$ un sous-groupe de $(G, \cdot)$. Les ensembles quotients $G / H$ et $H \backslash G$ sont équipotents. Plus précisément, l'application :
\[
\varphi : G / H \longrightarrow H \backslash G, \quad xH \longmapsto Hx^{-1}
\]
est bien définie et bijective.
\end{proposition}

\begin{demonstration}
\begin{itemize}
    \item \textbf{Vérifions que $\varphi$ est bien définie :}
    Soient $x, y \in G$ tels que $xH = yH$. Il s'agit de montrer que $Hx^{-1} = Hy^{-1}$.
    On a $x^{-1} y \in H$, donc il existe $h \in H$ tel que $y = xh$. Alors $y^{-1} = (xh)^{-1} = h^{-1} x^{-1}$, donc :
    \[
    Hy^{-1} = H(h^{-1} x^{-1}) = (H h^{-1}) x^{-1} = H x^{-1}
    \]
    (car $h^{-1} \in H$, donc $Hh^{-1} = H$). Ainsi $\varphi$ ne dépend pas du représentant choisi.
    
    \item \textbf{Bijectivité :} L'application réciproque de $\varphi$ est :
    \[
    \psi : H \backslash G \longrightarrow G / H, \quad Hx \longmapsto x^{-1} H
    \]
    (on vérifie de la même manière que $\psi$ est bien définie : si $Hx = Hy$, alors $x = hy$ avec $h \in H$, donc $x^{-1} = y^{-1} h^{-1}$, d'où $x^{-1} H = y^{-1} (h^{-1} H) = y^{-1} H$, et on a bien $\psi \circ \varphi = \mathrm{id}_{G/H}$ et $\varphi \circ \psi = \mathrm{id}_{H \backslash G}$).
\end{itemize}
\end{demonstration}

% ==============================================================================
% PAGE C12-44 & C12-45
% ==============================================================================

\begin{center}
  {\Large \textbf{\underline{Théorème de Lagrange}}}
\end{center}

\vspace{0.5em}

\begin{definition}
Soit $G$ un groupe fini. Le cardinal de $G$ est appelé l'\textbf{ordre} de $G$. On le note $|G|$.
\end{definition}

\vspace{0.5em}

\begin{theoreme}[Lagrange]
Soit $(G, \cdot)$ un groupe fini et $H$ un sous-groupe de $(G, \cdot)$. Alors l'ordre de $H$ divise l'ordre de $G$.
\end{theoreme}

\paragraph{Preuve :} Notons $G/H$ l'ensemble des classes à gauche modulo $H$ dans $G$.

C'est une partition de $G$, donc :
\[
G = \bigsqcup_{X \in G/H} X, \quad \text{la réunion étant disjointe.}
\]
Donc :
\[
\Card(G) = \sum_{X \in G/H} \Card(X)
\]
Or toute classe $X$ est équipotente à $H$ car de la forme $xH = \gamma_x(H)$ où $x \in G$.

En effet, $\gamma_x$, permutation de $G$, induit une bijection de $H$ sur $\gamma_x(H)$.

Donc $\Card(X) = \Card(H)$ pour tout $X \in G/H$.

Par suite :
\[
\Card(G) = \sum_{X \in G/H} \Card(X) = \Card(G/H)\Card(H) \qquad \blacksquare
\]

\begin{remarque}[NB]
La démonstration montre que :
\[
\Card(G/H) = \frac{\Card(G)}{\Card(H)}
\]
\end{remarque}

\vspace{0.5em}

\begin{exercice}[1]
\begin{enumerate}
  \item Soit $G$ un groupe fini d'ordre $n$ et $a \in G$ tel que $a \neq e$ et $a^2 = e$. Alors $n$ est pair.
  \item Soit $G$ un groupe fini d'ordre premier $p$ ($p$ est un nombre premier). Alors $G$ est monogène. En fait, pour tout $a \in G \setminus \{e\}$, $\langle a \rangle = G$.
\end{enumerate}
\end{exercice}

\paragraph{Preuve :}
\begin{enumerate}
  \item $H = \{e, a\}$ est un sous-groupe de $G$, d'ordre $2$. Par Lagrange, $|H| = 2$ divise $|G| = n$.
  \item Si $a \neq e$, le sous-groupe engendré par $a$, $\langle a \rangle$, est d'ordre $\ge 2$ car il contient $e$ et $a$.
  
  Par Lagrange, $|\langle a \rangle|$ divise $p = |G|$. Comme $p$ est premier et $|\langle a \rangle| > 1$, on a $|\langle a \rangle| = p = |G|$.
  
  Donc $G = \langle a \rangle$. \hfill $\blacksquare$
\end{enumerate}

% ==============================================================================
% PAGE C12-46
% ==============================================================================
\newpage

\begin{center}
  {\LARGE \textbf{\underline{IV. Morphisme de groupes}}}
\end{center}

\vspace{0.5em}

\subsection*{4.1. Définitions, exemples}

\begin{definition}
Soient $(G_1, \cdot)$ et $(G_2, \cdot)$ deux groupes (ce ne sont pas les mêmes lois même si les deux sont notées par un $\cdot$). On dit qu'une application $\varphi : G_1 \to G_2$ est un \textbf{morphisme de groupes} si :
\[
\forall x, y \in G_1, \quad \varphi(x \cdot y) = \varphi(x) \cdot \varphi(y)
\]
*(où $\cdot$ à gauche est la loi interne de $G_1$ et $\cdot$ à droite est la loi interne de $G_2$)*.
\end{definition}

\vspace{0.5em}

\begin{exemple}
\begin{enumerate}
  \item $\exp : (\C, +) \to (\C^*, \times)$, $z \mapsto e^z$ est un morphisme de groupes :
  \[
  \forall (z, z') \in \C^2, \quad e^{z + z'} = e^z \times e^{z'}
  \]
  \item $f : (\R, +) \to (\U, \times)$, $\theta \mapsto e^{i\theta}$ :
  \[
  \forall (\theta, \theta') \in \R^2, \quad f(\theta + \theta') = e^{i(\theta + \theta')} = e^{i\theta} \times e^{i\theta'}
  \]
  \item $\ln : (\R_+^*, \times) \to (\R, +)$ :
  \[
  \forall (x, y) \in (\R_+^*)^2, \quad \ln(xy) = \ln(x) + \ln(y)
  \]
  \item Soit $(G, \cdot)$ un groupe, et $x \in G$. L'application $\varphi_x : \Z \to G$, $n \mapsto x^n$ est un morphisme de groupes :
  \[
  \forall (m, n) \in \Z^2, \quad \varphi_x(m + n) = x^{m+n} = x^m \cdot x^n = \varphi_x(m) \cdot \varphi_x(n)
  \]
  En notation additive, $(G, +)$ étant commutatif : $\varphi_x : \Z \to G, \, n \mapsto n \cdot x$.
\end{enumerate}
\end{exemple}

% ==============================================================================
% PAGE C12-47 & C12-48
% ==============================================================================
\vspace{0.5em}

\begin{proposition}
Soient $(G_1, \cdot)$ et $(G_2, \cdot)$ deux groupes, de neutres respectifs $e_1$ et $e_2$. Alors :
\begin{enumerate}
  \item $\varphi(e_1) = e_2$
  \item $\forall x \in G_1, \quad \varphi(x^{-1}) = \varphi(x)^{-1}$
\end{enumerate}
\end{proposition}

\paragraph{Preuve :}
\begin{enumerate}
  \item On a $e_1 \cdot e_1 = e_1$ donc $\varphi(e_1 \cdot e_1) = \varphi(e_1)$, soit $\varphi(e_1) \cdot \varphi(e_1) = \varphi(e_1) = \varphi(e_1) \cdot e_2$.
  
  En simplifiant par $\varphi(e_1)$ à gauche dans le groupe $(G_2, \cdot)$, il vient :
  \[
  \varphi(e_1) = e_2
  \]
  \item Soit $x \in G_1$. Alors $x x^{-1} = e_1$ donc $\varphi(x x^{-1}) = \varphi(e_1) = e_2$, soit $\varphi(x) \cdot \varphi(x^{-1}) = e_2$.
  
  Comme on sait que $\varphi(x)$ est inversible dans $G_2$, on a nécessairement :
  \[
  \varphi(x^{-1}) = \varphi(x)^{-1} \qquad \blacksquare
  \]
\end{enumerate}

\begin{remarque}
On a aussi :
\begin{enumerate}
  \setcounter{enumi}{2}
  \item $\forall x \in G_1, \, \forall n \in \Z, \quad \varphi(x^n) = \varphi(x)^n$
\end{enumerate}
\begin{itemize}
  \item Vrai si $n = 0$ (car $x^0 = e_1$ dans $G_1$ et $y^0 = e_2$ dans $G_2$).
  \item Par récurrence dans $n \in \N$, c'est vrai pour tout $n \in \N$.
\end{itemize}
Puis, si $n \in \Z^-$ :
\[
\begin{aligned}
\varphi(x^n) = \varphi(x^{-|n|}) &= \varphi\left((x^{|n|})^{-1}\right) \\
&= \varphi(x^{|n|})^{-1} \\
&= \left(\varphi(x)^{|n|}\right)^{-1} \\
&= \varphi(x)^{-|n|} \\
&= \varphi(x)^n
\end{aligned}
\]
\end{remarque}

\begin{exemple}
$f : (\R, +) \to (\U, \times)$, $\theta \mapsto e^{i\theta}$ :
\[
\forall n \in \Z, \, \forall \theta \in \R, \quad f(n \cdot \theta) = f(\theta)^n \quad \text{c'est-à-dire} \quad e^{in\theta} = \left(e^{i\theta}\right)^n
\]
\end{exemple}

% ==============================================================================
% PAGE C12-48 (suite) & C12-49
% ==============================================================================
\vspace{0.5em}

\subsection*{4.2. Image directe / réciproque d'un sous-groupe par un morphisme}

\begin{proposition}
Soit $\varphi : (G_1, \cdot) \to (G_2, \cdot)$ un morphisme de groupes.
\begin{enumerate}
  \item Si $H$ est un sous-groupe de $(G_1, \cdot)$, alors $\varphi(H)$ est un sous-groupe de $(G_2, \cdot)$.
  \item Si $H$ est un sous-groupe de $(G_2, \cdot)$, alors $\varphi^{-1}(H)$ est un sous-groupe de $(G_1, \cdot)$.
\end{enumerate}
\end{proposition}

\paragraph{Preuve :}
\begin{enumerate}
  \item D'abord $\varphi(H) \subset G_2$.
  \begin{itemize}
    \item $e_1 \in H$ donc $\varphi(e_1) = e_2 \in \varphi(H)$.
    \item Soient $x_2, y_2 \in \varphi(H)$. Il existe donc $x_1, y_1 \in H$ tels que $x_2 = \varphi(x_1)$ et $y_2 = \varphi(y_1)$.
    
    Alors :
    \[
    x_2 y_2^{-1} = \varphi(x_1)\varphi(y_1)^{-1} = \varphi(x_1)\varphi(y_1^{-1}) = \varphi(x_1 y_1^{-1})
    \]
    Comme $H$ est un sous-groupe, $x_1 y_1^{-1} \in H$, donc $x_2 y_2^{-1} \in \varphi(H)$.
  \end{itemize}
  Donc $\varphi(H)$ est un sous-groupe de $(G_2, \cdot)$.

  \item Soit $H$ un sous-groupe de $(G_2, \cdot)$. Alors $\varphi^{-1}(H) = \{x \in G_1 \mid \varphi(x) \in H\}$.
  
  Déjà $\varphi^{-1}(H) \subset G_1$.
  \begin{itemize}
    \item $e_1 \in \varphi^{-1}(H)$ car $\varphi(e_1) = e_2 \in H$.
    \item Soient $x_1, y_1 \in \varphi^{-1}(H)$. Alors :
    \[
    \varphi(x_1 y_1^{-1}) = \varphi(x_1)\varphi(y_1)^{-1}
    \]
    Or $\varphi(x_1), \varphi(y_1) \in H$, qui est un sous-groupe, donc $\varphi(x_1) \cdot \varphi(y_1)^{-1} \in H$.
    
    Donc $x_1 y_1^{-1} \in \varphi^{-1}(H)$. \hfill $\blacksquare$
  \end{itemize}
\end{enumerate}

% ==============================================================================
% PAGE C12-49 (suite) & C12-50
% ==============================================================================
\vspace{0.5em}

\subsection*{4.3. Noyau et image d'un morphisme}

\begin{definition}
Soit $\varphi : (G_1, \cdot) \to (G_2, \cdot)$ un morphisme de groupes. On appelle \textbf{noyau} de $\varphi$, et on note $\Ker(\varphi)$ ou $\Ker \varphi$, l'ensemble $\varphi^{-1}(\{e_2\})$ où $e_2$ est le neutre de $(G_2, \cdot)$. Autrement dit :
\[
\Ker(\varphi) = \{x \in G_1 \mid \varphi(x) = e_2\}
\]
\end{definition}

\begin{exemple}
\begin{enumerate}
  \item $\exp : (\C, +) \to (\C^*, \times)$, $z \mapsto e^z$ :
  \[
  \Ker(\exp) = \{z \in \C \mid e^z = 1\} = 2i\pi\Z
  \]
  \item $f : (\R, +) \to (\U, \times)$, $\theta \mapsto e^{i\theta}$ :
  \[
  \Ker(f) = \{\theta \in \R \mid e^{i\theta} = 1\} = 2\pi\Z
  \]
  \item $\ln : (\R_+^*, \times) \to (\R, +)$, $x \mapsto \ln(x)$ :
  \[
  \Ker(\ln) = \{x \in \R_+^* \mid \ln x = 0\} = \{1\}
  \]
  \item $m \in \Z$, $\varphi_m : \Z \to \Z$, $x \mapsto m \cdot x$ :
  \begin{itemize}
    \item Si $m = 0$, $\Ker(\varphi_0) = \Z$.
    \item Si $m \neq 0$, $\Ker(\varphi_m) = \{x \in \Z \mid m \cdot x = 0\} = \{0\}$.
  \end{itemize}
\end{enumerate}
\end{exemple}

\begin{definition}
Soit $\varphi : (G_1, \cdot) \to (G_2, \cdot)$ un morphisme. On appelle \textbf{image} de $\varphi$, et on note $\Image(\varphi)$ ou $\Image \varphi$, l'ensemble $\varphi(G_1)$.
\end{definition}

\begin{proposition}
L'image et le noyau d'un morphisme de groupes $\varphi : (G_1, \cdot) \to (G_2, \cdot)$ sont des sous-groupes de $(G_2, \cdot)$ et $(G_1, \cdot)$, respectivement.
\end{proposition}

\paragraph{Preuve :}
\begin{itemize}
  \item $H = G_1$ est un sous-groupe de $(G_1, \cdot)$, donc $\varphi(H) = \varphi(G_1)$ est un sous-groupe de $(G_2, \cdot)$.
  \item $H = \{e_2\}$ est un sous-groupe de $(G_2, \cdot)$, donc $\varphi^{-1}(\{e_2\}) = \Ker(\varphi)$ est un sous-groupe de $(G_1, \cdot)$. \hfill $\blacksquare$
\end{itemize}

\begin{exemple}[Une preuve rapide de « $\U_n$ est un sous-groupe de $(\C^*, \times)$ »]
Soit $n \in \N^*$ et $\varphi_n : \C^* \to \C^*$, $z \mapsto z^n$.

C'est un morphisme de groupes car :
\[
\forall (z, z') \in \C^* \times \C^*, \quad (zz')^n = z^n z'^n
\]
Donc $\Ker(\varphi_n) = \{z \in \C^* \mid z^n = 1\} = \U_n$ est un sous-groupe de $(\C^*, \times)$.
\end{exemple}

% ==============================================================================
% PAGE C12-50 (bas) & C12-51 & C12-52
% ==============================================================================
\vspace{0.5em}

\begin{proposition}
Soit $\varphi : (G_1, \cdot) \to (G_2, \cdot)$ un morphisme de groupes, $a \in G_2$ et $x_0 \in G_1$ tel que $\varphi(x_0) = a$.

Alors l'ensemble des solutions $x \in G_1$ de l'équation $\varphi(x) = a$ est :
\[
x_0 \cdot \Ker(\varphi) = \Ker(\varphi) \cdot x_0
\]
\end{proposition}

\paragraph{Preuve :} $\forall x \in G_1$,
\[
\begin{aligned}
\varphi(x) = a &\iff \varphi(x) = \varphi(x_0) \\
&\iff \varphi(x)\varphi(x_0)^{-1} = e_2 \\
&\iff \varphi(x \cdot x_0^{-1}) = e_2 \\
&\iff x x_0^{-1} \in \Ker \varphi \\
&\iff \exists h \in \Ker \varphi, \quad x \cdot x_0^{-1} = h \\
&\iff \exists h \in \Ker \varphi, \quad x = h x_0 \\
&\iff x \in \Ker(\varphi) x_0
\end{aligned}
\]
On peut aussi écrire :
\[
\begin{aligned}
\varphi(x) = \varphi(x_0) &\iff \varphi(x_0)^{-1}\varphi(x) = e_2 \\
&\iff \varphi(x_0^{-1} x) = e_2 \\
&\iff x_0^{-1} x \in \Ker \varphi \\
&\iff \exists h \in \Ker \varphi, \quad x = x_0 h \\
&\iff x \in x_0 \cdot \Ker(\varphi)
\end{aligned}
\]

\paragraph{Lien avec les EDL :}
\[
\begin{aligned}
(E_f) &: \sum_{k=0}^p a_k(t) y^{(k)} = f(t) \qquad (t \in I) \\
(H) &: \sum_{k=0}^p a_k(t) y^{(k)} = 0
\end{aligned}
\]
Soient $G_1 : (\mathcal{D}^p(I, \R), +)$ et $G_2 : (\mathcal{F}(I, \R), +)$, et l'application :
\[
\varphi : \mathcal{D}^p(I, \R) \to \mathcal{F}(I, \R), \quad y \mapsto \sum_{k=0}^p a_k y^{(k)}
\]
On a $\forall y, z \in \mathcal{D}^p(I, \R)$, $\varphi(y + z) = \varphi(y) + \varphi(z)$.

$\Ker(\varphi) = \text{ensemble des solutions de l'équation homogène } (E_0) = \mathcal{S}_0$.

S'il existe $y_0 \in \mathcal{D}^p(I, \mathbb{K})$ tel que $\varphi(y_0) = f$, c'est-à-dire une solution particulière de $(E_f)$, alors l'ensemble des solutions de $(E_f)$ est :
\[
y_0 + \Ker \varphi = y_0 + \mathcal{S}_0
\]

\vspace{0.5em}

\begin{proposition}
Soit $\varphi : (G_1, \cdot) \to (G_2, \cdot)$ un morphisme de groupes.

Alors $\varphi$ est injectif si et seulement si $\Ker \varphi = \{e_1\}$.
\end{proposition}

\paragraph{Preuve :}
\begin{enumerate}
  \item Supposons $\varphi$ injectif, montrons que $\Ker \varphi = \{e_1\}$.
  
  D'abord, $e_1 \in \Ker \varphi$ car $\varphi(e_1) = e_2$.
  
  Réciproquement, soit $x \in \Ker \varphi$, c'est-à-dire $\varphi(x) = e_2$.
  
  Or $e_2 = \varphi(e_1)$, donc $\varphi(x) = \varphi(e_1)$ donc $x = e_1$ par injectivité de $\varphi$.
  
  Par double inclusion, on a bien $\Ker \varphi = \{e_1\}$.

  \item Supposons $\Ker \varphi = \{e_1\}$, montrons que $\varphi$ est injectif.
  
  Soient $x, y \in G_1$ tels que $\varphi(x) = \varphi(y)$.
  
  Alors $\varphi(x)\varphi(y)^{-1} = e_2$, soit $\varphi(x y^{-1}) = e_2$.
  
  Donc $x y^{-1} \in \Ker \varphi$. Or $\Ker \varphi = \{e_1\}$, donc $x y^{-1} = e_1$, donc $x = y$.
  
  Donc $\varphi$ est injective. \hfill $\blacksquare$
\end{enumerate}

% ==============================================================================
% PAGE C12-53 & C12-54 & C12-55
% ==============================================================================
\vspace{0.5em}

\begin{methode}
Pour prouver qu'un morphisme de groupes $\varphi : (G_1, \cdot) \to (G_2, \cdot)$ est injectif, il suffit de montrer que :
\[
\forall x \in G_1, \quad \varphi(x) = e_2 \implies x = e_1
\]
Il est inutile de rédiger la réciproque.
\end{methode}

\vspace{0.5em}

\subsection*{4.4. Isomorphisme de groupes}

Soit $E = \{0, 1\}$. On sait que $(\mathcal{P}(E), \Delta)$ *(où $\Delta$ = différence symétrique)* est un groupe d'ordre $4$. Le groupe produit $(\U_2 \times \U_2, \times)$ est également d'ordre $4$.

Leurs tables de Cayley respectives sont :

\begin{center}
\begin{tabular}{c|cccc}
$\Delta$ & $\emptyset$ & $\{0\}$ & $\{1\}$ & $E$ \\
\hline
$\emptyset$ & $\emptyset$ & $\{0\}$ & $\{1\}$ & $E$ \\
$\{0\}$ & $\{0\}$ & $\emptyset$ & $E$ & $\{1\}$ \\
$\{1\}$ & $\{1\}$ & $E$ & $\emptyset$ & $\{0\}$ \\
$E$ & $E$ & $\{1\}$ & $\{0\}$ & $\emptyset$ \\
\end{tabular}
\qquad \qquad
\begin{tabular}{c|cccc}
$\times$ & $(1, 1)$ & $(1, -1)$ & $(-1, 1)$ & $(-1, -1)$ \\
\hline
$(1, 1)$ & $(1, 1)$ & $(1, -1)$ & $(-1, 1)$ & $(-1, -1)$ \\
$(1, -1)$ & $(1, -1)$ & $(1, 1)$ & $(-1, -1)$ & $(1, 1)$ \\
$(-1, 1)$ & $(-1, 1)$ & $(-1, -1)$ & $(1, 1)$ & $(1, -1)$ \\
$(-1, -1)$ & $(-1, -1)$ & $(1, 1)$ & $(1, -1)$ & $(1, 1)$ \\
\end{tabular}
\end{center}

On remarque que ces $2$ tables coïncident si on identifie :
\[
\begin{aligned}
\emptyset &\text{ à } (1, 1) \\
\{0\} &\text{ à } (1, -1) \\
\{1\} &\text{ à } (-1, 1) \\
E &\text{ à } (-1, -1)
\end{aligned}
\]
En revanche, la table de Cayley de $(\U_4, \times)$ est essentiellement différente, car :
\[
\forall x \in \U_2 \times \U_2, \quad x^2 = (1, 1) \qquad (\forall X \in \mathcal{P}(E), \, X \Delta X = \emptyset)
\]
alors que $i^2 \neq 1$ (neutre de $\U_4$).

Pour cette raison, il n'existe aucune bijection de $(\U_2 \times \U_2, \times)$ vers $(\U_4, \times)$ qui identifie les $2$ tables.

On dira que $(\mathcal{P}(E), \Delta)$ et $(\U_2 \times \U_2, \times)$ sont \textbf{isomorphes}, mais $(\U_2 \times \U_2, \times)$ et $(\U_4, \times)$ ne le sont pas.

\vspace{0.5em}

\begin{definition}
Soient $(G_1, \cdot)$ et $(G_2, \cdot)$ deux groupes. On appelle \textbf{isomorphisme} ($\neq$ involution) $\varphi : (G_1, \cdot) \to (G_2, \cdot)$ tout morphisme bijectif.
\end{definition}

\begin{proposition}
Si $\varphi : (G_1, \cdot) \to (G_2, \cdot)$ est un isomorphisme, alors $\varphi^{-1} : (G_2, \cdot) \to (G_1, \cdot)$ est un morphisme (donc isomorphisme).
\end{proposition}

\paragraph{Preuve :} Notons $\varphi^{-1} : G_2 \to G_1$ l'application réciproque de la bijection $\varphi$.

Soient $x_2, y_2 \in G_2$. En notant $x_1 = \varphi^{-1}(x_2)$ et $y_1 = \varphi^{-1}(y_2)$, on a $x_1, y_1 \in G_1$ et $\varphi(x_1) = x_2$, $\varphi(y_1) = y_2$.

Alors :
\[
x_2 \cdot y_2 = \varphi(x_1) \cdot \varphi(y_1) = \varphi(x_1 \cdot y_1)
\]
Donc $x_1 y_1 = \varphi^{-1}(x_2 y_2)$, c'est-à-dire :
\[
\varphi^{-1}(x_2) \varphi^{-1}(y_2) = \varphi^{-1}(x_2 y_2)
\]
C'est vrai pour tout $(x_2, y_2) \in G_2^2$, donc $\varphi^{-1}$ est un morphisme de groupes. \hfill $\blacksquare$

\begin{remarque}
S'il existe un isomorphisme entre $(G_1, \cdot)$ et $(G_2, \cdot)$, on dit que $(G_1, \cdot)$ et $(G_2, \cdot)$ sont \textbf{isomorphes}.
\end{remarque}

\vspace{0.5em}

\begin{exemple}
\begin{enumerate}
  \item $\varphi : (\Z/2\Z, +) \to (\U_2, \times)$ définie par $\begin{cases} \dot{0} \mapsto 1 \\ \dot{1} \mapsto -1 \end{cases}$ est un isomorphisme.
  
  En effet, $\Z/2\Z = \{\dot{0}, \dot{1}\}$ avec $\dot{0} = 2\Z$ et $\dot{1} = 1 + 2\Z$ :
  \begin{center}
  \begin{tabular}{c|cc}
  $+$ & $\dot{0}$ & $\dot{1}$ \\
  \hline
  $\dot{0}$ & $\dot{0}$ & $\dot{1}$ \\
  $\dot{1}$ & $\dot{1}$ & $\dot{0}$ \\
  \end{tabular}
  \qquad \qquad
  \begin{tabular}{c|cc}
  $\times$ & $1$ & $-1$ \\
  \hline
  $1$ & $1$ & $-1$ \\
  $-1$ & $-1$ & $1$ \\
  \end{tabular}
  \end{center}
  
  \item Si $n \in \N^*$, on montrera que $(\Z/n\Z, +)$ et $(\U_n, \times)$ sont isomorphes.
  
  \item Plus généralement, si $E$ est de cardinal $n$, alors $(\mathcal{P}(E), \Delta)$ et $(\U_2^n, \times)$ (groupe-produit) sont isomorphes.
\end{enumerate}
\end{exemple}

\begin{exemple}[2 -- $(\U_2 \times \U_2, \times)$ et $(\U_4, \times)$ ne sont pas isomorphes]
Par l'absurde, supposons qu'il existe un isomorphisme $\varphi : \U_4 \to \U_2 \times \U_2$ *(on exhibe une propriété de groupe qui doit être invariante dans les deux groupes)*.

On a alors $\varphi(i)^2 = (1, 1)$ (car $\forall x \in \U_2 \times \U_2, \, x^2 = (1, 1)$), soit $\varphi(i^2) = (1, 1) = \text{le neutre de } \U_2 \times \U_2$.

Donc $i^2 = -1 \in \Ker \varphi$, donc $\Ker \varphi \neq \{1\}$. Donc $\varphi$ n'est pas injectif, *a fortiori* pas bijectif.
\end{exemple}

\section{V -- Structure quotient -- Applications}

\subsection{5.1. Relation d'équivalence compatible avec une loi interne}

\begin{definition}
Soit $(M, T)$ un magma et $\mathcal{R}$ une relation d'équivalence sur $M$.
\begin{itemize}
    \item On dit que $\mathcal{R}$ est \mot{compatible} avec $T$ si :
    \[
    \forall (x, x', y, y') \in M^4, \quad (x \, \mathcal{R} \, x' \ \text{et} \ y \, \mathcal{R} \, y') \implies (x \, T \, y) \, \mathcal{R} \, (x' \, T \, y')
    \]
    \item On dit que $\mathcal{R}$ est \mot{compatible à droite} avec $T$ si :
    \[
    \forall (x, y, z) \in M^3, \quad x \, \mathcal{R} \, y \implies (x \, T \, z) \, \mathcal{R} \, (y \, T \, z)
    \]
    \item On dit que $\mathcal{R}$ est \mot{compatible à gauche} avec $T$ si :
    \[
    \forall (x, y, z) \in M^3, \quad x \, \mathcal{R} \, y \implies (z \, T \, x) \, \mathcal{R} \, (z \, T \, y)
    \]
\end{itemize}
\end{definition}

\begin{proposition}
Une relation d'équivalence $\mathcal{R}$ sur $M$ est compatible avec $T$ si et seulement si elle est compatible à droite et à gauche avec $T$.
\end{proposition}

\begin{demonstration}
\begin{itemize}
    \item \textbf{Sens direct ($\implies$) :} Supposons $\mathcal{R}$ compatible avec $T$.
    Soit $(x, y) \in M^2$ tel que $x \, \mathcal{R} \, y$. Or, pour tout $z \in M$, $z \, \mathcal{R} \, z$ (par réflexivité).
    Donc, par compatibilité de $\mathcal{R}$ avec $T$, on a :
    \[
    (x \, T \, z) \, \mathcal{R} \, (y \, T \, z) \quad \text{et} \quad (z \, T \, x) \, \mathcal{R} \, (z \, T \, y)
    \]
    Ainsi, $\mathcal{R}$ est compatible à droite et à gauche avec $T$.

    \item \textbf{Réciproque ($\Longleftarrow$) :} Supposons $\mathcal{R}$ compatible à droite et à gauche avec $T$.
    Soit $(x, x', y, y') \in M^4$ tel que $x \, \mathcal{R} \, x'$ et $y \, \mathcal{R} \, y'$.
    \begin{itemize}
        \item Par compatibilité à droite : $(x \, T \, y) \, \mathcal{R} \, (x' \, T \, y)$.
        \item Par compatibilité à gauche : $(x' \, T \, y) \, \mathcal{R} \, (x' \, T \, y')$.
    \end{itemize}
    Donc, par transitivité de $\mathcal{R}$, $(x \, T \, y) \, \mathcal{R} \, (x' \, T \, y')$.
\end{itemize}
\end{demonstration}

\begin{theoreme}[et Définition]
Soit $(M, T)$ un magma et $\mathcal{R}$ une relation d'équivalence compatible avec $T$.
Soit la surjection canonique :
\[
s : M \longrightarrow M / \mathcal{R}, \quad x \longmapsto \overline{x}
\]
où $\overline{x}$ désigne la classe d'équivalence de $x$ \textit{(c'est l'application qui associe un représentant à sa classe d'équivalence)}.

Alors il existe une, et une seule loi de composition interne $\overline{T}$ sur $M / \mathcal{R}$ telle que :
\[
\forall (x, y) \in M^2, \quad s(x \, T \, y) = s(x) \, \overline{T} \, s(y)
\]
On dit que $\overline{T}$ est la \mot{loi quotient} de $T$ par $\mathcal{R}$.
\end{theoreme}

\begin{demonstration}
Procédons par analyse-synthèse :
\begin{itemize}
    \item \cas{Analyse :} On suppose que $\overline{T} : M / \mathcal{R} \times M / \mathcal{R} \to M / \mathcal{R}$ vérifie :
    \[
    \forall (x, y) \in M^2, \quad s(x \, T \, y) = s(x) \, \overline{T} \, s(y) \quad \text{c'est-à-dire} \quad \overline{x \, T \, y} = \overline{x} \, \overline{T} \, \overline{y}
    \]
    Soient $X, Y \in M / \mathcal{R}$. Choisissons $x \in X$ et $y \in Y$, c'est-à-dire tels que $X = \overline{x}$ et $Y = \overline{y}$.
    On doit avoir :
    \[
    X \, \overline{T} \, Y = \overline{x \, T \, y}
    \]
    Donc $\overline{T}$, sous réserve qu'elle soit bien définie, est unique.

    \item \cas{Synthèse :} Définissons $\overline{T} : M / \mathcal{R} \times M / \mathcal{R} \to M / \mathcal{R}$ par :
    \[
    (X, Y) \longmapsto \overline{x \, T \, y}
    \]
    où $x$ (resp. $y$) est n'importe quel élément de $X$ (resp. $Y$).

    Il faut vérifier que $\overline{T}$ est bien une application, ce qui revient à vérifier que pour tout $(X, Y) \in (M / \mathcal{R})^2$, le couple $(X, Y)$ a une seule image (indépendance vis-à-vis du choix des représentants).
    
    Soient $x, x' \in X$ et $y, y' \in Y$. On a $\overline{x} = \overline{x'} = X$ et $\overline{y} = \overline{y'} = Y$, donc $x \, \mathcal{R} \, x'$ et $y \, \mathcal{R} \, y'$.
    Par compatibilité de $\mathcal{R}$ avec $T$, on a $(x \, T \, y) \, \mathcal{R} \, (x' \, T \, y')$, donc :
    \[
    \overline{x \, T \, y} = \overline{x' \, T \, y'}
    \]
    Ainsi, $(X, Y)$ a une image unique. Par analyse-synthèse, il existe une unique loi $\overline{T} : M / \mathcal{R} \times M / \mathcal{R} \to M / \mathcal{R}$ telle que $\forall (x, y) \in M^2, \ s(x \, T \, y) = s(x) \, \overline{T} \, s(y)$.
\end{itemize}
\end{demonstration}

\begin{remarque}[NB]
Dans l'analyse, on montre que $\overline{T}$ a pour graphe :
\[
\big\{ ((X, Y), \, \overline{x \, T \, y}) \mid (X, Y) \in (M / \mathcal{R})^2, \ x \in X, \ y \in Y \big\}
\]
(ce n'était pas encore une fonction qu'on définissait, mais un graphe). Dans la synthèse, on vérifie que ce graphe est fonctionnel.
\end{remarque}

\begin{proposition}
Soit $(M, T)$ un magma, $\mathcal{R}$ une relation d'équivalence sur $M$ compatible avec $T$, et $\overline{T}$ la loi quotient sur $M / \mathcal{R}$.
\begin{enumerate}
    \item Si $T$ est associative, alors $\overline{T}$ aussi.
    \item Si $T$ est commutative, alors $\overline{T}$ aussi.
    \item Si $T$ possède un élément neutre $e$, alors $\overline{e}$ (la classe d'équivalence de $e$ modulo $\mathcal{R}$) est neutre pour $\overline{T}$.
    \item Si $(M, T)$ est un monoïde, alors $(M / \mathcal{R}, \overline{T})$ aussi.
\end{enumerate}
\end{proposition}

\begin{demonstration}
Notons $s : M \to M / \mathcal{R}, \ x \mapsto s(x) = \overline{x}$ la surjection canonique.
Les propriétés 1) à 3) découlent de la surjectivité de $s$ et de la propriété :
\[
\forall (x, y) \in M^2, \quad s(x \, T \, y) = s(x) \, \overline{T} \, s(y)
\]
\begin{itemize}
    \item \textbf{Preuve de 1) :} Soient $X, Y, Z \in M / \mathcal{R}$. Choisissons $x, y, z \in M$ tels que $X = s(x) = \overline{x}$, $Y = s(y) = \overline{y}$ et $Z = s(z) = \overline{z}$ (c'est-à-dire $x \in X$, $y \in Y$ et $z \in Z$). Alors :
    \begin{align*}
        (X \, \overline{T} \, Y) \, \overline{T} \, Z &= (s(x) \, \overline{T} \, s(y)) \, \overline{T} \, s(z) \\
                                                       &= s(x \, T \, y) \, \overline{T} \, s(z) \\
                                                       &= s((x \, T \, y) \, T \, z) \\
                                                       &= s(x \, T \, (y \, T \, z)) \quad (\text{par associativité de } T) \\
                                                       &= s(x) \, \overline{T} \, (s(y) \, \overline{T} \, s(z)) \\
                                                       &= X \, \overline{T} \, (Y \, \overline{T} \, Z)
    \end{align*}
    \item Les points 2) et 3) se démontrent de manière analogue.
    \item \textbf{Preuve de 4) :} Découle directement des points 1) et 3).
\end{itemize}
\end{demonstration}

\newpage
\subsection{5.2. Symétrisation d'un monoïde régulier}

\textbf{Objectif :} Donner une construction permettant d'obtenir le groupe $(\Z, +)$ à partir du monoïde $(\N, +)$.

Plus généralement, si $(M, +)$ est un monoïde abélien, on veut construire un groupe $(G, +)$ aussi petit que possible tel que $M \subset G$, et dont la loi prolonge celle de $M$.

Cela ne sera possible que si $(M, +)$ est un \mot{monoïde régulier}, c'est-à-dire :
\[
\forall m \in M, \ \forall (p, q) \in M^2, \quad m + p = m + q \implies p = q
\]

Soit $(M, +)$ un monoïde abélien régulier de neutre $0$. Munissons $M^2$ de la relation $\mathcal{R}$ définie par :
\[
\forall (m, n), (m', n') \in M^2, \quad (m, n) \, \mathcal{R} \, (m', n') \iff m + n' = m' + n
\]
*(L'idée sous-jacente est de s'affranchir de la soustraction en écrivant $m - n = m' - n'$ sous la forme $m + n' = m' + n$ pour pouvoir construire les symétriques).*

Munissons par ailleurs $M^2$ de la structure de \mot{monoïde produit} :
\[
\forall (m, n), (m', n') \in M^2, \quad (m, n) + (m', n') = (m + m', \, n + n')
\]
Il est facile de vérifier que cette addition produit est associative, commutative et possède $(0, 0)$ pour élément neutre.

\begin{theoreme}
Soit $(M, +)$ un monoïde abélien régulier. La relation $\mathcal{R}$ sur $M^2$ définie par :
\[
(m, n) \, \mathcal{R} \, (m', n') \iff m + n' = m' + n
\]
est une relation d'équivalence compatible avec la loi produit sur $M^2$.
\end{theoreme}

\begin{demonstration}
\begin{itemize}
    \item \textbf{Relation d'équivalence :} La réflexivité et la symétrie de $\mathcal{R}$ sont évidentes.
    
    \textit{Transitivité :} Soient $(m, n), (m', n'), (m'', n'') \in M^2$ tels que $(m, n) \, \mathcal{R} \, (m', n')$ et $(m', n') \, \mathcal{R} \, (m'', n'')$, c'est-à-dire :
    \[
    m + n' = m' + n \quad \text{et} \quad m' + n'' = m'' + n'
    \]
    Alors :
    \[
    m + n' + n'' = m' + n + n'' = (m' + n'') + n = m'' + n' + n
    \]
    Comme $n'$ est régulier dans $(M, +)$ (et la loi commutative), on peut simplifier par $n'$ et on obtient :
    \[
    m + n'' = m'' + n
    \]
    c'est-à-dire $(m, n) \, \mathcal{R} \, (m'', n'')$. Donc $\mathcal{R}$ est une relation d'équivalence.

    \item \textbf{Compatibilité avec la loi produit :} La loi produit sur $M^2$ étant commutative, il suffit de vérifier que $\mathcal{R}$ est compatible à gauche avec cette loi.
    
    Soient $(m, n), (m', n'), (m'', n'') \in M^2$ tels que $(m', n') \, \mathcal{R} \, (m'', n'')$, c'est-à-dire $m' + n'' = m'' + n'$.
    Vérifions que :
    \[
    \underbrace{(m, n) + (m', n')}_{(m + m', \, n + n')} \, \mathcal{R} \, \underbrace{(m, n) + (m'', n'')}_{(m + m'', \, n + n'')}
    \]
    On a, par associativité et commutativité de $+$ dans $M$ :
    \begin{align*}
        (m + m') + (n + n'') &= m + n + (m' + n'') \\
                             &= m + n + (m'' + n') \\
                             &= (m + m'') + (n + n')
    \end{align*}
    On a bien $\big((m, n) + (m', n')\big) \, \mathcal{R} \, \big((m, n) + (m'', n'')\big)$.
\end{itemize}
\end{demonstration}

D'après le paragraphe 5.1, l'ensemble quotient $M^2 / \mathcal{R}$ peut être muni de la loi quotient $\dot{+}$ telle que :
\[
\forall (m, n), (m', n') \in M^2, \quad [m, n] \, \dot{+} \, [m', n'] = [m + m', \, n + n']
\]
où $[m, n]$ désigne la classe d'équivalence de $(m, n) \in M^2$ modulo $\mathcal{R}$.

Notons $(G, \dot{+})$ le magma quotient $(M^2 / \mathcal{R}, \dot{+})$. On sait déjà, d'après le paragraphe 5.1, que $(G, \dot{+})$ est un monoïde commutatif dont le neutre est $[0, 0]$.

\begin{theoreme}
Soit $(M, +)$ un monoïde abélien régulier et $(G, \dot{+})$ le magma $(M^2 / \mathcal{R}, \dot{+})$ défini ci-dessus.
Alors :
\begin{enumerate}
    \item $(G, \dot{+})$ est un \mot{groupe abélien}.
    \item L'application :
    \[
    f : M \longrightarrow G, \quad m \longmapsto [m, 0]
    \]
    est un \mot{morphisme de monoïdes injectif}.
    \item Enfin, pour tout élément $g \in G$, il existe $(m, n) \in M^2$ tel que :
    \[
    g = f(m) - f(n)
    \]
\end{enumerate}
\end{theoreme}

\begin{demonstration}
On sait déjà que $(G, \dot{+})$ est un monoïde abélien de neutre $[0, 0]$.

\begin{itemize}
    \item \textbf{Symétrisabilité dans $G$ :} Soit $X \in G$ et $(m, n) \in M^2$ tel que $X = [m, n]$.
    Posons $X' = [n, m] \in G$. Alors :
    \[
    X \, \dot{+} \, X' = [m, n] \, \dot{+} \, [n, m] = [m + n, \, n + m] = [0, 0]
    \]
    (car $(m + n) + 0 = 0 + (n + m)$), et de même $X' \, \dot{+} \, X = [0, 0]$.
    Donc $(G, \dot{+})$ est un monoïde dont tout élément est symétrisable : c'est donc un groupe (abélien).

    \item \textbf{Propriétés de $f$ :} Considérons $f : M \to G, \ m \mapsto [m, 0]$.
    \begin{itemize}
        \item Elle vérifie $\forall m, n \in M, \ f(m + n) = [m + n, 0] = [m, 0] \, \dot{+} \, [n, 0] = f(m) \, \dot{+} \, f(n)$, ainsi que $f(0) = [0, 0]$ (condition nécessaire pour un morphisme de monoïdes). Donc $f$ est un morphisme de monoïdes.
        \item De plus, si $m, n \in M$ vérifient $f(m) = f(n)$, c'est-à-dire $[m, 0] = [n, 0]$, alors $(m, 0) \, \mathcal{R} \, (n, 0)$, soit $m + 0 = n + 0$, d'où $m = n$. Donc $f$ est injective.
    \end{itemize}
    Ainsi, $f(M)$, muni de la loi induite par $\dot{+}$, est un sous-monoïde du groupe $(G, \dot{+})$ isomorphe à $(M, +)$.

    \item \textbf{Décomposition de tout élément de $G$ :} On a montré que le symétrique de $[n, 0]$ est $-[n, 0] = [0, n]$.
    Donc, pour tout $g \in G$, en choisissant $(m, n) \in M^2$ tel que $g = [m, n]$, on a :
    \[
    g = [m, 0] \, \dot{+} \, [0, n] = [m, 0] - [n, 0] = f(m) - f(n)
    \]
    (en notant $[m, 0] - [n, 0]$ pour $[m, 0] \, \dot{+} \, (-[n, 0])$).
\end{itemize}
\end{demonstration}

\begin{remarque}
On peut identifier $M$ à $f(M)$, et donc noter simplement $m$ pour $f(m)$ si $m \in M$.
Par la propriété morphique de $f$, $m + n$ désigne aussi bien $f(m + n)$ que $f(m) \, \dot{+} \, f(n)$.
\end{remarque}

\begin{exemples}[Construction de $(\Z, +)$]
$(\N, +)$ est un monoïde abélien régulier. Par la construction décrite ci-dessus, on obtient le groupe $(\Z, +)$ (où la loi quotient $\dot{+}$ est notée $+$).

\textbf{Premières propriétés de $(\Z, +)$ :}
\begin{itemize}
    \item On a identifié $\N$ à $\Z^+ = \{ [m, 0] \mid m \in \N \}$.
    \item Tout $x \in \Z$ s'écrit $x = [m, n]$ où $(m, n) \in \N^2$ :
    \begin{itemize}
        \item Si $m \ge n$, alors $x = [m - n, 0] \in \Z^+ = \N$.
        \item Si $m \le n$, alors $x = [0, n - m] = -[n - m, 0]$, donc $-x \in \Z^+ = \N$.
    \end{itemize}
    Ainsi, $\Z = \N \cup (-\N)$.
    \item De plus, $\N \cap (-\N) = \{0\}$ car :
    \[
    [m, 0] = [0, n] \implies m + n = 0 \implies m = n = 0
    \]
    \item On note $\Z^+ = \N$ et $\Z^- = -\N$. L'application $\Z^+ \to \Z^-, \ x \mapsto -x$ est une bijection.
\end{itemize}

\begin{center}
\begin{tikzpicture}[scale=0.85]
  % Axes
  \draw[->, >=stealth, thick] (-4, 0) -- (5.5, 0);
  \draw[->, >=stealth, thick] (0, -0.5) -- (0, 5.2);

  % Graduations
  \foreach \x/\lbl in {-3/-3, -2/-2, -1/-1, 0/0, 1/1, 2/2, 3/3, 4.5/m} {
    \fill[bleu] (\x, 0) circle (2pt);
    \node[below=3pt, font=\small] at (\x, 0) {$\lbl$};
  }
  \node[below=3pt, font=\small] at (3.75, 0) {$\dots$};

  % Classes d'équivalence (droites obliques parallèles)
  \draw[bleu, thick] (-3, 0) -- (1.5, 4.5) node[right, font=\footnotesize\color{black}] {$[0, m]$};
  \draw[bleu, thick] (-2, 0) -- (2.1, 4.1) node[right, font=\footnotesize\color{black}] {$[0, 2]$};
  \draw[bleu, thick] (-1, 0) -- (2.7, 3.7) node[right, font=\footnotesize\color{black}] {$[0, 1]$};
  \draw[bordeaux, thick] (0, 0) -- (3.2, 3.2) node[right, font=\footnotesize\color{bordeaux}] {$[0, 0]$};
  \draw[bleu, thick] (1, 0) -- (3.8, 2.8) node[right, font=\footnotesize\color{black}] {$[1, 0]$};
  \draw[bleu, thick] (2, 0) -- (4.3, 2.3) node[right, font=\footnotesize\color{black}] {$[2, 0]$};
  \draw[bleu, thick] (3, 0) -- (4.8, 1.8) node[right, font=\footnotesize\color{black}] {$[3, 0]$};
  \draw[bleu, thick] (4.5, 0) -- (5.8, 1.3) node[right, font=\footnotesize\color{black}] {$[m, 0]$};

  % Note marginale
  \node[align=left, font=\footnotesize\itshape, text=brun] at (7.8, 2.5) {$\rightarrow$ Les entiers relatifs\\ sont des ensembles\\ (rien de surprenant :\\ les entiers naturels aussi).};
\end{tikzpicture}
\end{center}

\begin{itemize}
    \item Comme $(\Z, +)$ est un groupe abélien, la soustraction ne pose plus de problème : si $x, y \in \Z$, il existe un et un seul $d \in \Z$ tel que $x + d = y$, c'est $d = y - x$.
    \item Puisque $\Z = \N \cup (-\N)$ et $\N \cap (-\N) = \{0\}$, pour tout $x \in \Z$, un et un seul des éléments $x$ ou $-x$ appartient à $\N$. On le note $|x|$, appelé \mot{valeur absolue} de $x$. Pour tous $x, y \in \Z$, on a $|x + y| \le |x| + |y|$.
    \item On munit $\Z$ de la relation d'ordre :
    \[
    x \le y \iff y - x \in \N
    \]
    L'ordre est total puisqu'on a toujours $y - x \in \N$ ou $x - y \in \N$. On constate que cet ordre prolonge celui de $\N$.
    
    Enfin, l'ordre sur $\Z$ est compatible avec $+$ :
    \[
    \forall x, y, z \in \Z, \quad x \le y \implies x + z \le y + z
    \]
    et avec la multiplication par un entier $n > 0$ :
    \[
    \forall x, y \in \Z, \ \forall n \in \N^*, \quad x \le y \implies n \cdot x \le n \cdot y
    \]
    \item On munit $\Z$ d'une seconde loi de composition interne, la multiplication $\times$, telle que $\forall (x, y) \in \Z^2, \ x \times y = x \cdot y$.
    
    Cette loi prolonge la multiplication sur $\N$, et on obtient un \mot{anneau commutatif} $(\Z, +, \times)$. C'est un \mot{anneau intègre}, c'est-à-dire qu'il vérifie :
    \[
    \forall (x, y) \in \Z^2, \quad x \times y = 0 \implies x = 0 \ \text{ou} \ y = 0
    \]
\end{itemize}
\end{exemples}

\end{document}
